% Statements are numbered by hand, in the order they are stated.
\documentclass[11pt]{article}

\usepackage[margin=1in]{geometry}
\usepackage{lmodern}
\usepackage{microtype}
\usepackage{amsmath,amssymb,amsthm}
\usepackage{amsopn}
\usepackage{unicode-math}
\usepackage{amstext}
% A blank line in the source opens a line of space on the page.
% Section headings are adjusted so they do not pick up a second skip.
\usepackage[skip=\baselineskip]{parskip}
% Long inline formulas may break at a relation or a binary operator,
% instead of stretching the spaces on that line.
\relpenalty=100
\binoppenalty=200
% Reference links only. Section and statement cross-references stay black.
\usepackage[bookmarks=false,colorlinks,linkcolor=black,citecolor=black,urlcolor=blue]{hyperref}
\newcommand{\doilink}[1]{\href{https://doi.org/#1}{\nolinkurl{doi:#1}}}

\newenvironment{named}[1]{%
  \par\addvspace{\medskipamount}%
  \pagebreak[3]%
  \noindent{\bfseries #1}\space\ignorespaces
  \stmtoc{#1}%
}{\par\addvspace{\medskipamount}}

\newcommand{\topic}[1]{\par\noindent\textbf{#1}}

\setcounter{tocdepth}{2}
\ExplSyntaxOn
\prop_new:N \g_qn_toc_total_prop
\prop_new:N \g_qn_toc_seen_prop
\bool_new:N \g_qn_toc_appendix_bool
\AddToHook { cmd/appendix/before } { \bool_gset_true:N \g_qn_toc_appendix_bool }
\cs_new_protected:Npn \qnTocCount #1 #2
  { \prop_gput:Nnn \g_qn_toc_total_prop {#1} {#2} }
\cs_new_protected:Npn \qn_stmtoc:n #1
  {
    \tl_set:Nn \l_tmpa_tl {#1}
    \bool_set_false:N \l_tmpa_bool
    \tl_if_in:NnT \l_tmpa_tl { Outline } { \bool_set_true:N \l_tmpa_bool }
    \tl_if_in:NnT \l_tmpa_tl { Bauer } { \bool_set_true:N \l_tmpa_bool }
    \tl_if_in:NnT \l_tmpa_tl { Ridout } { \bool_set_true:N \l_tmpa_bool }
    \bool_if:NF \l_tmpa_bool
      {
        \regex_match:nVTF { \A (Conjecture|Theorem|Proposition)\b | conjecture } \l_tmpa_tl
          {
            \bool_if:NTF \g_qn_toc_appendix_bool
              {
                \regex_match:nVTF { \A Proposition \b } \l_tmpa_tl
                  { }
                  { \qn_stmtoc_write: }
              }
              { \qn_stmtoc_write: }
          }
          { }
      }
  }
\cs_new_protected:Npn \qn_stmtoc_write:
  {
    \tl_set_eq:NN \l_tmpb_tl \l_tmpa_tl
    \regex_replace_once:nnN { \.\Z } { } \l_tmpb_tl
    \int_set:Nn \l_tmpa_int { 1 }
    \prop_get:NVNT \g_qn_toc_seen_prop \l_tmpb_tl \l_tmpa_tl
      { \int_set:Nn \l_tmpa_int { \l_tmpa_tl + 1 } }
    \tl_set:Nx \l_tmpa_tl { \int_use:N \l_tmpa_int }
    \exp_args:NNVV \prop_gput:Nnn \g_qn_toc_seen_prop \l_tmpb_tl \l_tmpa_tl
    \use:e
      {
        \exp_not:N \qnTocWrite
        { \exp_not:V \l_tmpb_tl }
        { \exp_not:V \l_tmpa_tl }
      }
    \prop_get:NVNT \g_qn_toc_total_prop \l_tmpb_tl \l_tmpa_tl
      {
        \int_compare:nNnT { \l_tmpa_int } = { \l_tmpa_tl }
          {
            \phantomsection
            \exp_args:Nnx \addcontentsline {toc} {statement}
              { \exp_not:V \l_tmpb_tl }
          }
      }
  }
\newcommand{\stmtoc}[1]{ \qn_stmtoc:n {#1} }
\ExplSyntaxOff
\makeatletter
\newcommand{\qnTocWrite}[2]{%
  \immediate\write\@auxout{\string\qnTocCount{\unexpanded{#1}}{#2}}%
}
\newcommand*\l@statement[2]{%
  \ifnum \c@tocdepth > 1\relax
    \addpenalty\@secpenalty
    \begingroup
      \parindent \z@
      \leftskip 2.6em
      \rightskip \@pnumwidth
      \parfillskip -\@pnumwidth
      \leavevmode
      #1\nobreak\hfil\nobreak
      \hb@xt@\@pnumwidth{\hss #2}\par
    \endgroup
  \fi
}
\makeatother

\newcounter{stmtref}
% Conjectures and theorems share this counter.
% The headings are written out by hand, so a label would otherwise capture the section number.
\newcommand{\markstmt}[2]{%
  \setcounter{stmtref}{#1}%
  \addtocounter{stmtref}{-1}%
  \refstepcounter{stmtref}%
  \label{#2}%
}
% Lemmas, propositions, and corollaries share this counter.
\newcounter{lemmaref}
\newcommand{\marklemma}[2]{%
  \setcounter{lemmaref}{#1}%
  \addtocounter{lemmaref}{-1}%
  \refstepcounter{lemmaref}%
  \label{#2}%
}

\title{II. Cycles beyond $1$ for the $qn+1$ map: finiteness for each length, and an explicit bound on $q$ for a constant odd element}
\author{Trevor Hackman}
\date{}

\begin{document}
\maketitle

\begin{abstract}
% quotation indents every paragraph after the first and sets the paragraph skip to zero.
\setlength{\parskip}{\baselineskip}%
\setlength{\parindent}{0pt}%
\noindent
For the $qn+1$ map, for each length at least $2$, only finitely many odd integers $q$ have a cycle of that length.
The argument uses Ridout's theorem in the form that allows prime factors of the denominator outside a constant finite set, and it gives no bound on $q$.
If $\lvert 2^{r/\ell} - p/d \rvert > c\, d^{-1.3247}$ for each root $2^{r/\ell}$, for every odd positive integer $x$ and every large denominator $d = x\cdot 2^j$ or $d = \ell x\cdot 2^j$, then the bound on $q$ depends only on the length.

If a constant odd positive integer lies on a cycle with two or three distinct odd numbers, then $q$ lies below an explicit bound depending only on that integer.
For three odd numbers, once $q$ meets that bound, a lower bound and an upper bound cannot hold at the same time.
Bauer and Bennett (2018) give a restricted irrationality exponent of about $1.6149$, which lies just below the golden ratio, the threshold at which these bounds conflict.
Every smaller exponent is a finite check.
For two odd numbers the same conflict uses their exponent $1.465$ for the square root of $2$ (2002).
Evaluating their Theorem~1.5 at the exponent $1.617$, together with a search through $q < 2^{160{,}000}$ and a digit scan of $\sqrt{2}$, shows that every other cycle with two odd numbers has $q \ge 2^{1{,}071{,}999{,}971}$ and each odd element greater than $2^{17{,}631{,}663}$.
When the constant integer is $1$ and the cycle has three odd numbers, the bound is $q < 2^{17990}$, and the search in the companion paper already rules out every such cycle.
\end{abstract}

\begingroup
\setlength{\parskip}{0pt}
\small
\tableofcontents
\endgroup

\section{Introduction}
\label{sec:introduction}

This paper studies cycles of the $qn+1$ map for any starting value.
It is preceded by the narrower paper~\cite{hackman-i}, which studies cycles that pass through $1$.

Fix an odd integer $q \ge 3$.
For an odd positive integer $n$, let $v$ be the largest integer such that $2^v$ divides $qn+1$.
The map $C$ sends $n$ to
\[
C(n) = \frac{qn+1}{2^v}.
\]
The orbit of $n$ is the sequence $n$, $C(n)$, $C(C(n))$, \ldots.
When $q = 3$, $C$ is the odd-to-odd form of the Collatz map, also called the Syracuse function, and the Collatz conjecture, posed by Lothar Collatz in 1937 and still open, asserts that every orbit eventually reaches $1$.
For $q = 3$ the number $1$ is fixed, since $3 \cdot 1 + 1 = 4$.
For larger $q$, heuristic arguments suggest that most orbits grow without bound and that cycles are rare, although no orbit has been proved to have odd terms that tend to infinity.

The known cycles are these.
For each Mersenne number, the map has the cycle $1 \to 1$.
No other odd $q \ge 3$ has a cycle of length $1$.
For $q = 5$ there is $1 \to 3 \to 1$, and also $13 \to 33 \to 83 \to 13$ and $17 \to 43 \to 27 \to 17$.
For $q = 181$ there is $27 \to 611 \to 27$ and $35 \to 99 \to 35$.
Crandall recorded $27 \to 611$ from the identity $2^{15} - 181^2 = 7$.

\begin{named}{Conjecture 1.}
\markstmt{1}{conj:all}
These are all the cycles of $C$.
For every odd $q \ge 3$, every cycle is either $1 \to 1$ with $q$ a Mersenne number, or one of the five cycles at $q = 5$ and $q = 181$ listed above.
In particular, if $q \ge 7$ is neither $181$ nor a Mersenne number, then $C$ has no cycle.
\end{named}

There is then no cycle through $1$ other than $1 \to 1$ for a Mersenne number and $1 \to 3 \to 1$ for $q = 5$, which is the conjecture of Santos (2020), and for $q = 3$ there is no cycle other than $1 \to 1$, which is part of the Collatz conjecture.
Crandall (1978) conjectured that for every odd integer $q > 3$, some positive integer has an orbit that does not reach $1$~\cite{crandall}.
For $q > 3$ he expected an orbit whose odd terms tend to infinity.
The cycles at $q = 5$ and $q = 181$ do not have odd terms that tend to infinity.
Conjecture~1 does not assert that the odd terms of any orbit tend to infinity.
The searches through $1$ are in the companion paper~\cite{hackman-i}.
A search over cycles with no element held constant, and a search at a constant $q$, are in Appendix~\ref{app:computations}.
Section~\ref{sec:all-cycles} gives a heuristic for the expected number of cycles.

\begin{named}{Theorem 2.}
\markstmt{2}{thm:finite}
For each $\ell \ge 2$, only finitely many odd $q$ have a cycle of $C$ of length $\ell$.
\end{named}

Theorem~2 is ineffective: it rests on the general form of Ridout's theorem~\cite{ridout} in Section~\ref{sec:ineffective}, which allows prime factors of the denominator outside a constant finite set, and gives no bound on $q$, and no bound depending on a constant starting value.
(A finiteness result is called effective if its proof gives an explicit bound on any exception, so that a finite computation can find them all, and ineffective if it gives no such bound.)
The case in which the cycle visits $1$ is Theorem~3 of the companion paper, and that proof uses only the $S$-unit form of Ridout's theorem.
The proof of Theorem~2 is in Section~\ref{sec:ineffective}.

\begin{named}{Outline of Theorem 3.}
Fix a length at least $2$ and an odd starting value.
Suppose each root of $2$ that this length produces has a restricted irrationality exponent below the golden ratio with respect to the denominators built from that starting value and a power of $2$.
Then any $q$ for which the starting value lies on a cycle of this length lies below a bound computed from that exponent.
For length $2$ the restricted irrationality exponent is $1.465$, from Corollary~1.4 of Bauer and Bennett (2002).
For length $3$ it is about $1.6149$, from Bauer and Bennett (2018).
The precise statement is in Section~\ref{sec:ql}.
\end{named}

Proposition~\ref{prop:exponent} gives, for each length, a bound on $q$ that does not depend on a starting value.
The hypothesis is a restricted irrationality exponent at most $1.3247$ for every root $2^{r/\ell}$, in every denominator $x\cdot 2^j$ and $\ell x\cdot 2^j$ with $x$ an odd positive integer.

For $x = 1$ and length $3$, Appendix~\ref{app:computations} gives $Q_3(1) < 2^{17990}$.
Any such cycle has one candidate $q$ for each total number of factors of $2$ removed, so it is enough to check those totals below $53970$.
The companion paper checks every such total up to $200{,}000$ and finds no cycle.
That search includes the totals below $53970$, and Theorem~1 of~\cite{hackman-i} records that no such cycle exists.
Length $2$ through $1$ is Proposition~A.1 of~\cite{hackman-i}: the only such $q$ is $5$.

For a general $x$ the bound does not rule a cycle out.
At $q = 181$ the cycles $27 \to 611 \to 27$ and $35 \to 99 \to 35$ both occur.
Appendix~\ref{app:computations} records $181 < Q_2(27)$ and $181 < Q_2(35)$.

Section~\ref{sec:length-2} takes up cycles of two odd numbers with no element held constant in advance.
Every such cycle other than $1 \to 3 \to 1$, $27 \to 611 \to 27$ and $35 \to 99 \to 35$ has $q \ge 2^{1{,}071{,}999{,}971}$ and every odd number on it greater than $2^{17{,}631{,}663}$.

Length $4$ with an odd total number of removed factors of $2$ remains open, as in the companion paper: the roots are $2^{1/4}$ and $2^{3/4}$, and a constant odd part does not supply an exponent below the golden ratio for them.
Corollary~\ref{cor:longer} covers the longer classes whose roots are only $\sqrt{2}$, or only the cube roots of $2$.

\section{Definitions and notation}
\label{sec:definitions}

Throughout, $q \ge 3$ is an odd integer and $C$ is the map defined in Section~\ref{sec:introduction}.

\topic{Mersenne numbers.}
Let $e \ge 2$ be an integer.
A Mersenne number is the integer $2^e - 1$.

\topic{The $2$-adic valuation.}
For a nonzero integer $n$, $v_2(n)$ is the largest integer $v$ such that $2^v$ divides $n$.

\topic{Cycles.}
A cycle of $C$ of length $\ell \ge 2$ consists of $\ell$ distinct odd positive integers.
If those integers are $x_0, x_1, \ldots, x_{\ell-1}$, then $C(x_{i-1}) = x_i$ and $C(x_{\ell-1}) = x_0$, where the indices are read modulo $\ell$.

\topic{Halvings.}
The halving counts $k_1, \ldots, k_\ell$ of the cycle are defined by
\[
q x_{i-1} + 1 = 2^{k_i} x_i
\qquad (1 \le i \le \ell,\ x_\ell = x_0),
\]
so $k_i = v_2(q \cdot x_{i-1} + 1)$.
The total number of halvings is $h = k_1 + \cdots + k_\ell$, and $S_i = k_1 + \cdots + k_i$.

\topic{The restricted irrationality exponent.}
Let $\theta$ be an irrational real number and let $\mathcal{D}$ be a set of positive integers.
A real number $\mu$ is a restricted irrationality exponent of $\theta$ with respect to $\mathcal{D}$ when there exist an explicit constant real number $c > 0$ and an explicit constant positive integer $d_0$ such that, for every integer $p$ and every positive integer $d \in \mathcal{D}$ with $d \ge d_0$,
\[
\left|\theta - \frac{p}{d}\right| > c \cdot d^{-\mu}.
\]
Write $\mathcal{D}_\ell$ for the set of denominators $2^j$ and $\ell \cdot 2^j$, where $j$ is a nonnegative integer.

\topic{The golden ratio.}
The golden ratio is $\varphi = (1 + \sqrt{5})/2$.

\topic{The plastic number.}
The plastic number is the real root of $t^3 - t - 1 = 0$.
The plastic number equals
\[
\sqrt[3]{\frac{9 + \sqrt{69}}{18}} + \sqrt[3]{\frac{9 - \sqrt{69}}{18}}.
\]

\topic{The constants $\mu_3$, $\gamma_3$, and $c_3(s')$.}
Let $N = 16166467234$.
The real number $\mu_3$ is $3\log_2 N/63$.
The positive number $\gamma_3$ is $1 - \mu_3 + 1/\mu_3$.
For a positive integer $s'$, the constant $c_3(s')$ is $2^{-5}\, N^{-79/63}/s'$.

\section{Ineffective finiteness}
\label{sec:ineffective}

\begin{named}{Ridout's theorem (1958).}
Let $\alpha$ be an algebraic irrational, let $S$ be a finite set of primes, and let $\varepsilon > 0$.
Only finitely many reduced fractions $p/Q$ with $Q > 0$ have every prime factor of $Q$ in $S$ and satisfy
\[
\left| \alpha - \frac{p}{Q} \right| < Q^{-1-\varepsilon}.
\]
The same theorem has the general form: writing $\lvert n \rvert_p = p^{-v_p(n)}$ for a nonzero integer $n$, only finitely many rationals $P/Q$ in lowest terms satisfy
\[
Q \bigl| Q\alpha - P \bigr| \prod_{p \in S} \lvert Q \rvert_p < \max\bigl(\lvert P \rvert, Q\bigr)^{-\varepsilon}.
\]
\end{named}

The first statement is the $S$-unit form.
If every prime factor of $Q$ lies in $S$, then $\prod_{p \in S} \lvert Q \rvert_p = 1/Q$ and the left side of the general form equals $\lvert Q\alpha - P \rvert$.
Whenever $\lvert P \rvert \le Q$, this is the first form.
The general form allows prime factors of $Q$ outside $S$.
Such a prime does not appear in the product, and it contributes its full size to the height $\max(\lvert P \rvert, Q)$.
Neither form gives a bound on the exceptional fractions.
Since $\alpha$ is irrational, the finitely many fractions that meet the general form for a constant $\alpha$ lie at a positive distance from $\alpha$, and these distances have a positive minimum.
This is why Theorem~\ref{thm:finite}, which uses the general form in place of a restricted irrationality exponent, gives no bound on $q$.

\begin{named}{Lemma 1 (the window for any cycle).}
\marklemma{1}{lem:any-window}
Let $a = x_0 \to x_1 \to \cdots \to x_{\ell-1} \to a$ be a cycle of length $\ell \ge 2$, with $a$ its smallest element, and let $\sigma = \sum_{i=0}^{\ell-1} 1/x_i$.
Then
\[
2^h = \prod_{i=0}^{\ell-1}\left(q + \frac{1}{x_i}\right),
\]
so $q^\ell < 2^h < (q+1)^\ell$, and in particular $\ell \nmid h$ and $q = \lfloor 2^{h/\ell} \rfloor$.
Writing $2^{h/\ell} = q + \delta$,
\[
e^{-\ell/q} \le \frac{\ell\,\delta}{\sigma} \le e^{\ell/q}.
\]
\end{named}

\begin{proof}
Multiplying $2^{k_i} x_i = q\, x_{i-1} + 1$ around the cycle gives the product.
Each factor lies in $(q,\, q+1]$, with equality to $q+1$ only for an element equal to $1$.
At most one element is $1$, and $\ell \ge 2$, so the product is strictly larger than $q^\ell$ and strictly smaller than $(q+1)^\ell$.
A power of $2$ strictly between consecutive $\ell$-th powers is not an $\ell$-th power, so $\ell \nmid h$, and $q = \lfloor 2^{h/\ell} \rfloor$.
In particular $0 < \delta < 1$.

The binomial theorem gives
\[
\ell q^{\ell-1}\delta \le (q+\delta)^\ell - q^\ell \le \ell q^{\ell-1}\delta\Bigl(1 + \frac{\delta}{q}\Bigr)^{\ell-1}.
\]
Dividing the binomial expansion by $\ell q^{\ell-1}\delta$, the coefficient of $(\delta/q)^k$ is $\binom{\ell}{k+1}/\ell = \binom{\ell-1}{k}/(k+1) \le \binom{\ell-1}{k}$, and $\sum_{k \ge 0}\binom{\ell-1}{k}(\delta/q)^k = (1 + \delta/q)^{\ell-1}$.
Since $\delta < 1$, this factor is at most $e^{(\ell-1)/q} \le e^{\ell/q}$.

The same difference equals $\sum_{j \ge 1} e_j q^{\ell-j}$, where $e_j$ is the $j$-th elementary symmetric function of the reciprocals $1/x_i$.
Thus $e_1 = \sigma$, and the right side is at least $q^{\ell-1}\sigma$.
In the expansion of $\sigma^j$, each product of $j$ distinct reciprocals occurs $j!$ times, so $e_j \le \sigma^j/j!$.
Hence the right side is at most
\[
\sum_{j \ge 1} \frac{\sigma^j}{j!}\, q^{\ell-j} = q^{\ell-1} \sum_{m \ge 0} \frac{\sigma^{m+1}}{(m+1)!\, q^m} \le q^{\ell-1}\sigma\, e^{\sigma/q},
\]
since $(m+1)! \ge m!$.
Every element is at least $1$, so $\sigma \le \ell$ and $e^{\sigma/q} \le e^{\ell/q}$.
Comparing the two sides gives $\ell\delta \le \sigma\, e^{\ell/q}$ and $\ell\delta\, e^{\ell/q} \ge \sigma$.
\end{proof}

\begin{named}{Lemma 2 (the length-$3$ gap).}
\marklemma{2}{lem:length-3-gap}
Let the cycle have length $3$.
Let $D = 2^h - q^3$.
Then $D \ge q^{4/3}$, except when $q = 5$ and $h = 7$.
In that case $D = 3$.
\end{named}

\begin{proof}
Lemma~\ref{lem:any-window} gives $q^3 < 2^h$, so $D > 0$ and $q^3 \ne 2^h$.
Bauer and Bennett (2018)~\cite{bauer-bennett-2018} prove that if $x$ and $n$ are integers with $x^3 \ne 2^n$ and $(x, n) \ne (5, 7)$, then $\lvert x^3 - 2^n \rvert \ge x^{4/3}$.
Taking $x = q$ and $n = h$ gives $D \ge q^{4/3}$, except when $(q, h) = (5, 7)$.
In that case $2^7 - 5^3 = 128 - 125 = 3$.
\end{proof}

For a cycle of length $2$, Section~\ref{sec:length-2} proves a lower bound on $2^h - q^2$ that requires $h$ large.

For a cycle of length at least $4$, Bauer and Bennett (2018) does not give a lower bound $2^h - q^\ell \ge q^{\theta}$ for a constant positive real number $\theta$.
Appendix~\ref{app:denominator} records the constraints on $D$ that hold at every length.

\begin{named}{Lemma 3 (three approximations).}
\marklemma{3}{lem:three-approx}
In the setting of Lemma~\ref{lem:any-window}, write $h = \ell m + r$ with $0 < r < \ell$.
Let $\theta_i = 2^{ir/\ell}$ for $1 \le i \le \ell-1$, let $\theta = \theta_1$, and let $y = \min(x_1, \ldots, x_{\ell-1})$.
Suppose $q \ge 4\ell$.
Then $2^m > q/2$, and:

\begin{enumerate}
\item $\lvert \theta - q/2^m \rvert = \delta/2^m$.
\item For $1 \le i \le \ell-1$, the rational number $F_i = x_i/(a \cdot 2^{im - S_i})$ satisfies $\lvert \theta_i - F_i \rvert \le 2i\theta_i/(aq)$.
\item The rational number $W = (\ell aq + 1)/(\ell a \cdot 2^m)$ satisfies $\lvert \theta - W \rvert \le 2(1/y + 2\ell/q)/q$.
\end{enumerate}
\end{named}

\begin{proof}
Since $2^m \theta = q + \delta$ and $\theta = 2^{r/\ell} < 2$, one has $2^m = (q+\delta)/\theta > q/2$, and part~1 follows.

The relation $q\, x_{i-1} + 1 = 2^{k_i} x_i$ gives $x_i\, 2^{S_i} = a \prod_{u < i}(q + 1/x_u)$ by induction on $i$, so
\[
F_i = \frac{1}{2^{im}}\prod_{u < i}\Bigl(q + \frac{1}{x_u}\Bigr), \qquad \theta_i = \frac{(q+\delta)^i}{2^{im}}.
\]
Lemma~\ref{lem:any-window} gives $\delta \le (\sigma/\ell)\, e^{\ell/q}$.
Every element is at least $a$, so $\sigma \le \ell/a$ and $\delta \le c/a$ with $c = e^{\ell/q} \le e^{1/4}$.
Each factor $q + 1/x_u$ and the number $q+\delta$ therefore lie in $[q,\, q + c/a]$, and both numerators lie in $[q^i,\, (q + c/a)^i]$.
Their difference is at most
\[
\Bigl(q + \frac{c}{a}\Bigr)^i - q^i = q^i\Bigl(\Bigl(1 + \frac{c}{aq}\Bigr)^i - 1\Bigr) \le \frac{ic}{a}\, q^{i-1}\, e^{ic/(aq)},
\]
where $e^z - 1 \le z\, e^z$.
Dividing by $2^{im} = (q+\delta)^i/\theta_i \ge q^i/\theta_i$ gives $\lvert \theta_i - F_i \rvert \le (ic\theta_i/(aq))\, e^{ic/(aq)}$.
Here $i \le \ell - 1$ and $aq \ge 4\ell$, so $ic/(aq) < c/4$.
For $0 < z \le 1/3$ and $n \ge 2$, one has $n! \ge 2^{n-1}$, so $z^n/n! \le z^n/2^{n-1}$ and $\sum_{n \ge 2} z^n/2^{n-1} = z^2/(2-z)$.
Thus $e^{1/4} < 1 + 1/4 + (1/16)/(7/4) = 9/7 < 4/3$ and $e^{1/3} < 1 + 1/3 + 1/15 = 7/5$.
Since $c \le e^{1/4} < 4/3$, one has $c/4 < 1/3$ and $c\, e^{c/4} < (4/3) \cdot (7/5) = 28/15 < 2$.
Thus $\lvert \theta_i - F_i \rvert \le 2i\theta_i/(aq)$.

For part~3,
\[
\bigl\lvert \theta - W \bigr\rvert = \frac{1}{2^m}\Bigl\lvert \delta - \frac{1}{\ell a} \Bigr\rvert.
\]
Write $\delta - 1/(\ell a) = (\sigma - 1/a)/\ell + (\delta - \sigma/\ell)$.
The other elements are at least $y$, so $0 \le \sigma - 1/a \le (\ell-1)/y$ and $(\sigma - 1/a)/\ell \le 1/y$.
From the two bounds in Lemma~\ref{lem:any-window}, with $z = \ell/q$,
\[
\frac{\sigma}{\ell} - \delta \le \frac{\sigma}{\ell}\bigl(1 - e^{-z}\bigr), \qquad \delta - \frac{\sigma}{\ell} \le \frac{\sigma}{\ell}\bigl(e^z - 1\bigr).
\]
The inequality $1 - e^{-z} \le e^z - 1 \le z\, e^z$ and $\sigma \le \ell$ give $\lvert \delta - \sigma/\ell \rvert \le (\ell/q)\, e^{\ell/q}$.
Since $q \ge 4\ell$, this is at most $(\ell/q)\, e^{1/4} < 2\ell/q$.
Therefore $\lvert \delta - 1/(\ell a) \rvert \le 1/y + 2\ell/q$.
Dividing by $2^m > q/2$ gives part~3.
\end{proof}

\begin{named}{Theorem 2.}
For each $\ell \ge 2$, only finitely many odd $q$ have a cycle of $C$ of length $\ell$.
\end{named}

\begin{proof}[Proof of Theorem~\ref{thm:finite}]
Fix $\ell \ge 2$, and set $\varepsilon = 1/(10\ell)$.
Let $S$ consist of $2$ and the prime factors of $\ell$, and let $\Theta$ be the set of real numbers $2^{t/\ell}$ with $1 \le t \le \ell(\ell-1)$ and $\ell \nmid t$.
Each element of $\Theta$ is an algebraic irrational.
Apply the general form of Ridout's theorem to each $\theta' \in \Theta$, with this $S$ and $\varepsilon$.
Only finitely many fractions meet the inequality, and none equals $\theta'$, so there is an $\eta > 0$ such that every such fraction satisfies $\lvert \theta' - P/d \rvert \ge \eta$.
The value of $\eta$ is not provided by Ridout's theorem.
The odd $q < 4\ell$ are finite in number, so take an odd $q \ge 4\ell$ with a cycle of length $\ell$, in the notation of Lemmas~\ref{lem:any-window} and~\ref{lem:three-approx}.
Then $\theta \in \Theta$.
Each $\theta_i$ either lies in $\Theta$ or is a power of $2$.

Consider the fraction $q/2^m$, which is in lowest terms.
Its left side in the general form equals $2^m \cdot \delta \cdot 2^{-m} = \delta$, since $\lvert 2^m \rvert_p = 1$ for every odd prime $p$.
Also $\max(q,\, 2^m) < q+1$, because $2^m = (q+\delta)/\theta < (q+1)/\theta < q+1$.
If $\delta < (q+1)^{-\varepsilon}$, then $\delta < \max(q,\, 2^m)^{-\varepsilon}$, so Ridout's inequality applies and $\delta/2^m \ge \eta$.
Combined with $2^m > q/2$ this yields $q < 2/\eta$.

Now assume $\delta \ge (q+1)^{-\varepsilon}$.
Lemma~\ref{lem:any-window} gives $\delta \le e^{\ell/q}/a \le e^{1/4}/a$, so $a < 2(q+1)^\varepsilon$.

Suppose first that $y \le q^{1-2\varepsilon}$, and choose $i$ with $x_i = y$.
If $\ell$ divides $ir$, then $\theta_i = 2^s$ for some integer $s \ge 1$.
This power is not equal to $F_i$.
Otherwise $x_i = a \cdot 2^{im - S_i + s}$; the right side is an odd positive integer only if the exponent of $2$ vanishes, which forces $x_i = a$, contradicting distinctness.
Let $e = \max(im - S_i,\, 0)$.
In lowest terms the denominator of $F_i$ divides $a \cdot 2^e$, so $\lvert \theta_i - F_i \rvert \ge 1/(a \cdot 2^e)$.
Lemma~\ref{lem:three-approx} then gives $2^e \ge q/(2i\theta_i)$.
If $e = 0$, then $q \le 2i\theta_i$.
Here $i \le \ell-1$ and $\theta_i < 2^{\ell-1}$, so $q < (\ell-1)\, 2^\ell$.
If $e \ge 1$, the same lemma and $q \ge 4\ell$ give $\lvert \theta_i - F_i \rvert \le 2(\ell-1)\theta_i/(4\ell) < \theta_i/2$, hence $F_i > \theta_i/2$.
Therefore
\[
y = F_i \cdot a \cdot 2^e > \frac{\theta_i}{2} \cdot a \cdot \frac{q}{2i\theta_i} = \frac{aq}{4i} \ge \frac{q}{4(\ell-1)}.
\]
Once $q > \bigl(4(\ell-1)\bigr)^{5\ell}$, the right side exceeds $q^{1-2\varepsilon}$, because $1/(2\varepsilon) = 5\ell$.
This contradicts $y \le q^{1-2\varepsilon}$.

If instead $\ell$ does not divide $ir$, then $\theta_i \in \Theta$.
Write $F_i = P/d$ in lowest terms.
The odd part of $d$ divides $a$, so $d \cdot \lvert d \rvert_2 \le a$.
Every factor $\lvert d \rvert_p$ is at most $1$ and $2 \in S$, so the left side of the general form is at most
\[
d^2 \cdot \lvert \theta_i - F_i \rvert \cdot \lvert d \rvert_2 \le \frac{2i\theta_i\, d}{q}.
\]
As above, $\lvert \theta_i - F_i \rvert < \theta_i/2$, so $F_i > \theta_i/2 > 1/2$ and $F_i < (3/2)\theta_i < 2^\ell$.
If $im \ge S_i$, the fraction before reducing is $y/(a \cdot 2^{im-S_i})$, whose denominator equals $y/F_i < 2y$, and the numerator is $y$.
If $im < S_i$, that denominator is $a$, so $d \le a \le y$, and the numerator is $F_i d < 2^\ell y$.
In both cases $d < 2y$ and $\max(\lvert P \rvert,\, d) < 2^\ell y$.
Also $2i\theta_i < (\ell-1)\, 2^\ell$, so the left side is less than $(\ell-1)\, 2^{\ell+1} y/q$.
The right side is greater than $(2^\ell y)^{-\varepsilon}$.
The comparison
\[
(\ell-1)\, 2^{\ell+1}\, y^{1+\varepsilon} < q \cdot 2^{-\ell\varepsilon}
\]
therefore implies Ridout's inequality.
It holds for every admissible $y$ once
\[
q^{\varepsilon + 2\varepsilon^2} > (\ell-1)\, 2^{\ell + 1 + \ell\varepsilon},
\]
because $y^{1+\varepsilon} \le q^{(1-2\varepsilon)(1+\varepsilon)} = q^{1 - \varepsilon - 2\varepsilon^2}$.
The right side depends only on $\ell$.
For every larger $q$, Ridout's theorem gives $\lvert \theta_i - F_i \rvert \ge \eta$.
Lemma~\ref{lem:three-approx} gives $2i\theta_i/(aq) \ge \eta$, and $2i\theta_i < (\ell-1)\, 2^\ell$ with $a \ge 1$, so $q < (\ell-1)\, 2^\ell/\eta$.

Suppose instead that $y > q^{1-2\varepsilon}$.
The numerator $N = \ell aq + 1$ is congruent to $1$ modulo $\ell a$, so $\gcd(N,\, \ell a) = 1$.
Any common factor of $N$ and $\ell a \cdot 2^m$ is therefore a power of $2$.
Let $v = v_2(N)$ and $e = \max(m - v,\, 0)$.
In lowest terms, $W = P/d$ with $d = \ell a \cdot 2^e$ and $P \le N$.
The product formula gives $\prod_p \lvert \ell \rvert_p = 1/\ell$, and every prime factor of $\ell$ lies in $S$, so $\prod_{p \in S} \lvert \ell \rvert_p = 1/\ell$.
Since $\lvert a \rvert_p \le 1$,
\[
\prod_{p \in S} \lvert d \rvert_p \le \frac{2^{-e}}{\ell}.
\]
The left side of the general form is therefore at most
\[
d^2 \cdot \lvert \theta - W \rvert \cdot \frac{2^{-e}}{\ell} = \ell a^2\, 2^e \cdot \lvert \theta - W \rvert \le \ell a^2\, 2^m \cdot \lvert \theta - W \rvert.
\]
Lemma~\ref{lem:three-approx} and $2^m < q+1$ bound this by $2\ell a^2(1 + 1/q)(1/y + 2\ell/q)$.
Insert $1/y < q^{-1+2\varepsilon}$ and $a < 2(q+1)^\varepsilon$.
For $q \ge 4\ell$ one has $1 + 1/q < 2$, $1 + 2\ell\, q^{-2\varepsilon} \le 1 + 2\ell$, and $(q+1)^{2\varepsilon} < (2q)^{2\varepsilon}$, so the left side is less than
\[
16\ell(1+2\ell)\, 2^{2\varepsilon}\, q^{-1+4\varepsilon}.
\]
Since $P \le \ell aq + 1$ and $d \le \ell a \cdot 2^m < \ell a(q+1)$, the height satisfies $\max(\lvert P \rvert,\, d) < \ell a(q+1) + 1 < 4\ell(q+1)^{1+\varepsilon}$, so the right side is greater than $(4\ell)^{-\varepsilon}(q+1)^{-\varepsilon(1+\varepsilon)} > (4\ell)^{-\varepsilon}(2q)^{-\varepsilon(1+\varepsilon)}$.
The quotient of these two bounds is
\[
16\ell(1+2\ell)\, 2^{2\varepsilon}(4\ell)^\varepsilon\, 2^{\varepsilon(1+\varepsilon)}\, q^{-1+5\varepsilon+\varepsilon^2}.
\]
The exponent equals $-1 + 5/(10\ell) + 1/(100\ell^2) \le -1 + 1/4 + 1/400 < -1/2$.
The quotient tends to $0$ as $q \to \infty$, so Ridout's inequality holds for all $q$ above a bound depending only on $\ell$.
Then $\lvert \theta - W \rvert \ge \eta$.
Lemma~\ref{lem:three-approx} and $1/y \le 1$ give $\lvert \theta - W \rvert \le 2(1+2\ell)/q$, hence $q \le (2+4\ell)/\eta$.

In every case $q$ is at most a number depending only on $\ell$ and $\eta$.
\end{proof}

The argument does not produce a bound on $q$, and it does not produce a bound depending on a constant odd start.
Fix an integer $V \ge 2$.
Each length $\ell$ with $2 \le \ell \le V$ contributes finitely many odd $q$, and for each such pair the inequalities $q^\ell < 2^h < (q+1)^\ell$ leave finitely many $h$ and finitely many compositions of $h$ into $\ell$ positive parts.
Thus only finitely many cycles have length at most $V$.

\emph{Theorem~2 fails for $\ell = 1$.}
A Mersenne number $q$ gives the cycle $1 \to 1$, and the same identity becomes $2^h = q + 1$.
The constant $\eta$ is not explicit, so neither is the bound on $q$.
An effective bound would require an effective inequality of the same shape for every denominator whose odd part is at most a constant times $q^\varepsilon$.

\section{The integer $Q_\ell(x)$}
\label{sec:ql}

\begin{named}{Theorem 3.}
\markstmt{3}{thm:bound}
Let $\ell \ge 2$ be an integer and let $x$ be an odd positive integer.
Let $\mu$ be a constant real number with $1 < \mu < \varphi$.
Let $c$ and $c_*$ be constant positive real numbers, and let $j_0$ be a positive integer.
Suppose that for every $r \in \{1, \ldots, \ell-1\}$, every integer $p$, and every integer $j \ge j_0$,
\[
\Bigl| 2^{r/\ell} - \frac{p}{x\, 2^{j}} \Bigr| > c \cdot 2^{-\mu j}
\]
and
\[
\Bigl| 2^{r/\ell} - \frac{p}{\ell x\, 2^{j}} \Bigr| > c_* \cdot 2^{-\mu j}.
\]
Let $B$ be the positive real number $\max\bigl(\ell-1,\ 2(\ell-2) + 1/x\bigr)$.
Let $\gamma$ be the positive real number $1 - \mu + 1/\mu$, and let $A$ and $Q_*$ be the positive real numbers
\[
A = \frac{x}{2}\, 2^{1/\ell}
\Biggl( \frac{c}{2^{(\ell-1)/\ell}\, B} \Biggr)^{1/\mu},
\qquad
Q_* = \Biggl( \frac{(\ell-1)/(\ell A) + \ell}{c_*\, 2^{1-\mu}} \Biggr)^{1/\gamma}.
\]
For each such $r$, let $d_r$ be the positive real number equal to the distance from $x\, 2^{r/\ell}$ to the nearest integer.
For each such $r$, each integer $j$ with $0 \le j < j_0$, and each odd positive integer $y < 3x \cdot 2^{j_0}$, let $\rho$ be the nonnegative real number
\[
\rho = \Biggl| \frac{y}{x\, 2^{j}} \cdot 2^{-r/\ell} - 1 \Biggr|.
\]
Whenever $\rho > 0$, the corresponding value $B/\rho$ is a positive real number, and so is $B \cdot 2^{r/\ell} \cdot x / d_r$.
Let $Q_0$ be the maximum of those finitely many positive real numbers.
If $\ell$ is composite, that maximum also includes the positive real number $Bx$.
Let $Q_e$ be the least positive integer such that every integer $q \ge Q_e$ satisfies $q \ge 4\ell$ and $e^{\ell/q} \le q^{1-1/\mu}$.
If $\ell$ is prime, let $Q_A$ be the positive integer $1$.
If $\ell$ is composite, let $Q_A$ be the least positive integer such that every integer $q \ge Q_A$ satisfies $q/(2B) > A\, q^{1/\mu}$.
Let $Q_\ell(x;\, \mu, c, c_*, j_0)$ be the positive real number equal to the maximum of $2^{j_0}$, of $Q_0$, of $Q_*$, of $Q_e$, and of $Q_A$.
If $x$ lies on a cycle of length $\ell$, then $q < Q_\ell(x;\, \mu, c, c_*, j_0)$.
\end{named}

\begin{proof}[Proof of Theorem~\ref{thm:bound}]
Let $x = x_0 \to x_1 \to \cdots \to x_{\ell-1} \to x$ be a cycle of length $\ell$, with each $x_i$ an odd positive integer.
Lemma~\ref{lem:any-window} gives $q = \lfloor 2^{h/\ell} \rfloor$ and $\ell \nmid h$.
Write $2^{h/\ell} = q + \delta$ with $\delta$ the fractional part, so $0 < \delta < 1$.
Let $m$ be the nonnegative integer and $r$ the integer with $0 < r < \ell$ and $h = \ell m + r$.
If $q < 2^{j_0}$ or $q < Q_e$ or $q < Q_A$, there is nothing to prove.
Assume $q \ge 2^{j_0}$, $q \ge Q_e$ and $q \ge Q_A$.
Then $q \ge 4\ell$ and $e^{\ell/q} \le q^{1-1/\mu}$.
If $\ell$ is composite, then $q/(2B) > A\, q^{1/\mu}$.
Since $1/x \le 1$, one has $B \le 2\ell - 3$, and therefore $2B < 4\ell \le q$.

For each integer $i$ with $1 \le i \le \ell-1$, multiplying $q x_u + 1 = 2^{k_{u+1}} x_{u+1}$ for $0 \le u < i$ gives
\[
x_i\, 2^{S_i} = x \prod_{u < i}\Bigl(q + \frac{1}{x_u}\Bigr).
\]
Define the positive real numbers $\varepsilon_i$ by $\varepsilon_1 = 1/(qx)$ and
\[
\varepsilon_i = \varepsilon_{i-1} + \frac{1+\varepsilon_{i-1}}{q\, x_{i-1}}
\qquad (i \ge 2).
\]
Then $x_i\, 2^{S_i} = x\, q^i(1+\varepsilon_i)$: the case $i = 1$ is $qx + 1 = xq\bigl(1 + 1/(qx)\bigr)$, and $q x_{i-1} + 1 = 2^{k_i} x_i$ multiplies by $q + 1/x_{i-1}$.
The bound $q \ge 4\ell$ gives $B/q < 1$ and $\varepsilon_1 = 1/(qx) \le 1/q$.
If $\varepsilon_{i-1} \le \bigl(1/x + 2(i-2)\bigr)/q$, then that upper bound is at most $B/q < 1$, and $x_{i-1} \ge 1$, so
\[
\varepsilon_i \le \varepsilon_{i-1} + \frac{2}{q} \le \frac{1/x + 2(i-1)}{q}.
\]
Thus $\varepsilon_i \le B/q$ for every $i \le \ell-1$.
Let $n_i$ be the integer $ih - \ell S_i$, so $n_i \equiv ir \pmod{\ell}$.
Let $j_i$ be the integer and $r_i$ the integer with $n_i = \ell j_i + r_i$ and $0 \le r_i < \ell$.
Let $\rho_i$ be the real number given by $x_i = x\, 2^{n_i/\ell}(1+\rho_i)$.
Then
\[
1+\rho_i = \Bigl(\frac{q}{q+\delta}\Bigr)^i (1+\varepsilon_i).
\]
Since $0 < \delta < 1$, one has $\bigl(q/(q+\delta)\bigr)^i < 1$ and $\bigl(q/(q+\delta)\bigr)^i \ge \bigl(q/(q+1)\bigr)^i \ge 1 - i/(q+1) > 1 - i/q$.
Also $1+\varepsilon_i \ge 1$, so $\rho_i < \varepsilon_i \le B/q$ and $\rho_i > -i/q$.
Thus $|\rho_i| < B/q < 1/2$.

If $\ell$ is prime, then $r_i \ne 0$: here $1 \le i \le \ell-1$ and $1 \le r \le \ell-1$, so $ir \not\equiv 0 \pmod{\ell}$.
If $r_i = 0$, then $\ell$ is composite and $x_i = x\, 2^{j_i}(1+\rho_i)$.
If $j_i = 0$, then $x_i$ and $x$ are distinct odd positive integers, so $|x_i - x| \ge 2$.
Hence $|\rho_i| \ge 2/x$ and $q < Bx/2$.
If $j_i \le -1$, write $t = -j_i$.
Then $x_i\, 2^{t} = x(1+\rho_i)$.
The integers $x_i\, 2^{t}$ and $x$ are distinct, because $x$ is odd, so $|x_i\, 2^{t} - x| \ge 1$.
Hence $|\rho_i| \ge 1/x$ and $q < Bx$.
In both cases $Bx$ is one of the numbers whose maximum is $Q_0$, so $q < Q_0$.
If $j_i \ge j_0$, then $j_i \ge 1$, so $x\, 2^{j_i}$ is even and $x_i$ is odd.
Hence $x\, 2^{j_i}\, |\rho_i| \ge 1$.
Combined with $|\rho_i| < B/q$ and $1+\rho_i > 1/2$, this gives $x_i > q/(2B) > A\, q^{1/\mu}$.

Now suppose $r_i \in \{1, \ldots, \ell-1\}$.
If $j_i < 0$, write $t = -j_i$, so $t$ is a positive integer and $x_i\, 2^{t}/x = 2^{r_i/\ell}(1+\rho_i)$.
The left side is an integer divided by $x$, so
\[
|\rho_i| \ge 2^{-r_i/\ell}\, \frac{d_{r_i}}{x}.
\]
Combined with $|\rho_i| < B/q$, this gives $q < B \cdot 2^{r_i/\ell} \cdot x / d_{r_i}$.
That value is one of the numbers whose maximum is $Q_0$, and $q < Q_0$.
If $0 \le j_i < j_0$, then $x_i < 3x \cdot 2^{j_i}$, because $2^{r_i/\ell} < 2$ and $1+\rho_i < 3/2$.
The value $B/|\rho_i|$ is one of the numbers whose maximum is $Q_0$, and $|\rho_i| < B/q$ gives $q < Q_0$.
If $j_i \ge j_0$, the hypothesis gives $2^{r_i/\ell}\, |\rho_i| > c \cdot 2^{-\mu j_i}$.
Since $2^{r_i/\ell} \le 2^{(\ell-1)/\ell}$ and $2^{r_i/\ell}(1+\rho_i) > 2^{1/\ell}\cdot(1/2)$,
\[
2^{j_i} > \Biggl( \frac{c\, q}{2^{(\ell-1)/\ell}\, B} \Biggr)^{1/\mu}
\]
and $x_i > A\, q^{1/\mu}$.

Every $x_i$ with $i \ge 1$ is therefore greater than $A\, q^{1/\mu}$, or else $q < Q_0$.
In the remaining case, the positive real number $\sigma = \sum_{i=0}^{\ell-1} 1/x_i$ satisfies $\sigma - 1/x \le (\ell-1)\, A^{-1} q^{-1/\mu}$.
Let $z$ be the positive real number $\ell/q$.
Lemma~\ref{lem:any-window} gives $e^{-z} \le \ell\delta/\sigma \le e^{z}$, so $|\ell\delta/\sigma - 1| \le e^{z} - 1$.
The series for $e^{z}$ gives $e^{z} - 1 \le z\, e^{z}$.
Since $\sigma \le \ell$,
\[
\bigl|\delta - \sigma/\ell\bigr| \le \frac{\ell}{q}\, e^{\ell/q}.
\]
Hence
\[
\bigl|\delta - 1/(\ell x)\bigr| \le \frac{\ell-1}{\ell A}\, q^{-1/\mu} + \frac{\ell}{q}\, e^{\ell/q}.
\]
The hypothesis $q \ge Q_e$ makes the second summand at most $\ell\, q^{-1/\mu}$, so
\[
\bigl|\delta - 1/(\ell x)\bigr| \le \Biggl( \frac{\ell-1}{\ell A} + \ell \Biggr) q^{-1/\mu}.
\]
Since $q \ge 2^{j_0}$ and $2^{h/\ell} > q$, one has $m \ge j_0$.
Also $2^m < (q+1)\, 2^{-r/\ell} < 2q$.
The window hypothesis, with denominator $\ell x \cdot 2^m$, gives
\[
\bigl|\delta - 1/(\ell x)\bigr|
= 2^m \Bigl| 2^{r/\ell} - \frac{\ell x q + 1}{\ell x \cdot 2^m} \Bigr|
> c_* \cdot 2^{(1-\mu)m}
> c_*\, 2^{1-\mu}\, q^{1-\mu}.
\]
Since $\gamma > 0$, these two bounds on $|\delta - 1/(\ell x)|$ conflict for every integer $q \ge Q_*$.
The lower bound is strict, so the conflict includes $q = Q_*$.
Thus $q < Q_*$, and in every case $q < Q_\ell(x;\, \mu, c, c_*, j_0)$.
\end{proof}

For $\ell = 2$ one has $B = 1$ and $Q_A = 1$.
Section~\ref{sec:q2} uses the larger factor $1+1/x$ in place of $B$, and records that the resulting maximum equals $Q_2(x)$.

For $\ell = 3$ one has $B = 2 + 1/x$ and $Q_A = 1$.
Section~\ref{sec:q3} uses the smaller factor $1+1/x$.
The proof of Corollary~\ref{cor:q3} shows that this factor bounds both $|\rho|$ of the other two elements once $q \ge 30$, and $30$ lies below $Q_3(x)$.

\begin{named}{Proposition 4.}
\marklemma{4}{prop:exponent}
Let $\ell \ge 2$ be an integer.
Let $\mu$ be a constant real number with $1 < \mu \le 1.3247$.
Let $c$ be an explicit constant real number with $c > 0$, and let $d_0$ be an explicit constant positive integer.
Suppose that for every integer $r$ with $1 \le r \le \ell - 1$, every odd positive integer $x$, every nonnegative integer $j$, and every integer $p$, each positive integer $d \in \{x \cdot 2^j,\, \ell x \cdot 2^j\}$ with $d \ge d_0$ satisfies
\[
\Bigl| 2^{r/\ell} - \frac{p}{d} \Bigr| > c \cdot d^{-\mu}.
\]
Then there is a positive real number $Q$, computed from $\ell$, $\mu$, $c$ and $d_0$, such that every cycle of length $\ell$ has $q < Q$.
\end{named}

\begin{proof}
Let $\gamma = 1 - \mu + 1/\mu$ and let $\alpha = \mu(\mu - 1)/\gamma$.
Since $1.3247 < 3/2 < \varphi$, one has $\mu < \varphi$, so $\mu^2 - \mu - 1 < 0$ and $\gamma > 0$.
For a real number $t > 1$ the derivative of $t^3 - t - 1$ is $3t^2 - 1 > 0$, so $t^3 - t - 1$ increases with $t$.
At $t = 1.3247 = 13247/10000$,
\[
\Bigl(\frac{13247}{10000}\Bigr)^3 - \frac{13247}{10000} - 1 = -\frac{76579777}{10^{12}} < 0.
\]
Thus $\mu^3 - \mu - 1 < 0$.
Since $\mu > 0$, this is the same inequality as $\mu^2(\mu - 1) < 1 + \mu - \mu^2$.
Dividing by $\mu$ gives $\mu(\mu - 1) < (1 + \mu - \mu^2)/\mu$.
The right side equals $\gamma$, so $\mu(\mu - 1) < \gamma$.
Since $\gamma > 0$, one has $\alpha < 1$.
Let $\beta = \mu(\mu - 1)$.
For $t > 1$ the real number $t(t - 1)$ increases with $t$, so
\[
\beta \le \frac{13247}{10000} \cdot \frac{3247}{10000} = \frac{43013009}{100000000} < \frac12.
\]

Let $x$ be the smallest element of a cycle of length $\ell$, and write $h = \ell m + r$ with $0 < r < \ell$, as in Lemma~\ref{lem:any-window}.
Then $2^m \cdot 2^{r/\ell} = q + \delta$, so $\lvert 2^{r/\ell} - q/2^m \rvert = \delta/2^m$.
If $2^m \ge d_0$, the hypothesis at the denominator $2^m$ gives $\delta/2^m > c \cdot 2^{-\mu m}$, so $\delta > c \cdot 2^{(1-\mu)m}$.
Lemma~\ref{lem:any-window} gives $\delta \le e^{\ell/q}/x$.
The proof of Lemma~\ref{lem:three-approx} records $e^{1/4} < 4/3$.
If $q \ge 4\ell$ and $2^m \ge d_0$, then $e^{\ell/q} \le e^{1/4}$ and
\[
x < \frac{4}{3c} \cdot 2^{(\mu-1)m}.
\]
Also $q > 2^{m + r/\ell} - 1$.
For $m \ge 1$ this is greater than $2^{m-1}$, so $\log_2 q > m - 1$.

If $2^m < d_0$, then $q < 2 \cdot 2^m < 2d_0$.
Let $Q_e$ be the least positive integer such that every integer $q \ge Q_e$ satisfies $q \ge 4\ell$ and $e^{\ell/q} \le q^{1 - 1/\mu}$, as in Theorem~\ref{thm:bound}.
Let $B_+ = 2\ell - 3$.
For $\ell = 2$ let $B_- = 1$, and for $\ell \ge 3$ let $B_- = \max(\ell - 1,\, 2(\ell - 2))$.
In Theorem~\ref{thm:bound} the positive real number $B$ satisfies $B_- \le B \le B_+$.
With the constant in that theorem equal to $c x^{-\mu}$, the factor $x$ cancels $x^{-\mu}$ under the exponent $1/\mu$, and
\[
A = \frac12 \cdot 2^{1/\ell}
\Biggl( \frac{c}{2^{(\ell-1)/\ell}\, B} \Biggr)^{1/\mu}.
\]
Set
\[
A_+ = \frac12 \cdot 2^{1/\ell}
\Biggl( \frac{c}{2^{(\ell-1)/\ell}\, B_-} \Biggr)^{1/\mu},
\qquad
A_- = \frac12 \cdot 2^{1/\ell}
\Biggl( \frac{c}{2^{(\ell-1)/\ell}\, B_+} \Biggr)^{1/\mu}.
\]
Then $A_- \le A \le A_+$.
If $\ell$ is prime, let $Q_A^\dagger = 1$.
If $\ell$ is composite, let $Q_A^\dagger$ be the least positive integer such that every integer $q \ge Q_A^\dagger$ satisfies $q/(2B_+) > A_+ \, q^{1/\mu}$.
If $q/(2B_+) > A_+ \, q^{1/\mu}$, then $q/(2B) > A \, q^{1/\mu}$.
Thus the positive integer $Q_A$ in Theorem~\ref{thm:bound} is at most $Q_A^\dagger$.

If $d_0 \le 1$, let $Q_1 = 1$.
If $d_0 > 1$, then for each odd positive integer $z < d_0$ let $j_0(z)$ be the least integer $j \ge 1$ with $z \cdot 2^j \ge d_0$.
Theorem~\ref{thm:bound}, with constants $c \, z^{-\mu}$ and $c \, (\ell z)^{-\mu}$ and with this $j_0(z)$, gives $q < Q_\ell(z;\, \mu,\, c z^{-\mu},\, c(\ell z)^{-\mu},\, j_0(z))$ whenever $z$ lies on a cycle of length $\ell$.
Let $Q_1$ be the maximum of those finitely many values.

Now suppose $x \ge d_0$, $2^m \ge d_0$ and $q \ge \max(Q_e,\, Q_A^\dagger)$.
Then $q \ge 4\ell$.
For every integer $j \ge 1$ the denominators $x \cdot 2^j$ and $\ell x \cdot 2^j$ are at least $d_0$, so
\[
\Bigl| 2^{r/\ell} - \frac{p}{x \cdot 2^j} \Bigr| > c \, x^{-\mu} \cdot 2^{-\mu j}
\]
and
\[
\Bigl| 2^{r/\ell} - \frac{p}{\ell x \cdot 2^j} \Bigr| > c \, (\ell x)^{-\mu} \cdot 2^{-\mu j}.
\]
Theorem~\ref{thm:bound}, with these constants and with $j_0 = 1$, gives $q < Q_\ell(x;\, \mu,\, c x^{-\mu},\, c(\ell x)^{-\mu},\, 1)$.
In that maximum, $2^{j_0} = 2$.
Since $A \ge A_-$ and $c_* = c(\ell x)^{-\mu}$,
\[
Q_* \le \Biggl( \frac{K \, (\ell x)^{\mu}}{c \cdot 2^{1-\mu}} \Biggr)^{1/\gamma},
\qquad
K = \frac{\ell-1}{\ell A_-} + \ell.
\]
Let
\[
\Lambda = \frac1\gamma \log_2\Biggl( \frac{K \, \ell^{\mu}}{c \cdot 2^{1-\mu}} \Biggr) + \frac{\mu}{\gamma} \log_2\Biggl( \frac{4}{3c} \Biggr).
\]
From $x < (4/(3c)) \cdot 2^{(\mu-1)m}$ one has $\log_2 Q_* < \Lambda + \alpha m$.

For $0 \le j < 1$ the only such integer is $j = 0$.
The hypothesis at denominator $x$ gives, for the distance $d_r$ from $x \cdot 2^{r/\ell}$ to the nearest integer, $d_r > c \, x^{1-\mu}$.
The value $B \cdot 2^{r/\ell} \cdot x / d_r$ in the maximum $Q_0$ is then less than $(2B_+/c) \, x^{\mu}$.
For an odd positive integer $y$ and for $\rho = \lvert y x^{-1} 2^{-r/\ell} - 1 \rvert$, the same hypothesis gives $B/\rho < (2B_+/c) \, x^{\mu}$ whenever $\rho > 0$.
If $\ell$ is composite, the value $Bx$ in that maximum satisfies $Bx \le B_+ x \le B_+ x^{\mu}$.
Let $K_0 = \max\bigl(B_+,\, 2B_+/c\bigr)$.
Thus $Q_0 \le K_0 x^{\mu}$, and
\[
\log_2(K_0 x^{\mu}) < \log_2 K_0 + \mu \log_2\Biggl( \frac{4}{3c} \Biggr) + \beta m.
\]

Let $m_0$ be the least integer $m \ge \max\bigl(1,\, \lceil \log_2 d_0 \rceil\bigr)$ such that
\[
\Lambda + \alpha m < m - 1
\]
and
\[
\log_2 K_0 + \mu \log_2\Biggl( \frac{4}{3c} \Biggr) + \beta m < m - 1.
\]
Such an integer exists because $\alpha < 1$ and $\beta < 1/2$.
If $m \ge m_0$, both $\log_2 Q_*$ and $\log_2(K_0 x^{\mu})$ are less than $m - 1$, hence less than $\log_2 q$.
Then $Q_* < q$ and $K_0 x^{\mu} < q$.
Together with $q \ge Q_e$, $q \ge Q_A^\dagger$ and $q > 2$, this contradicts $q < Q_\ell(x;\, \mu,\, c x^{-\mu},\, c(\ell x)^{-\mu},\, 1)$.

Therefore $m < m_0$, or $q < 2d_0$, or $q < Q_1$, or $q < \max(Q_e,\, Q_A^\dagger)$.
In the first of these, $q < 2^{m+1} \le 2^{m_0}$.
Let $Q$ be the maximum of $2d_0$, of $Q_1$, of $Q_e$, of $Q_A^\dagger$ and of $2^{m_0}$.
Every cycle of length $\ell$ has $q < Q$.
\end{proof}

The positive root of $t^3 - t - 1 = 0$ is the plastic number, and the plastic number is greater than $1.3247$.
For a constant real number $\mu$ below the plastic number one has $\mu^3 - \mu - 1 < 0$, and the same argument produces a bound on $q$.

\section{Cube roots with a constant odd part}
\label{sec:measure}

Bauer and Bennett~\cite{bauer-bennett-2018} prove that $\mu_3$ is a restricted irrationality exponent of $2^{1/3}$ and of $2^{2/3}$ when the denominator has odd part $1$ or $3$.
The inequality used here is the one obtained in the proof of their Theorem~2.1 from their inequality~(5.3), Lemma~4.1 and Proposition~5.1, before they round $\mu_3$ to $1.62$.
The companion paper records it for odd part $1$ or $3$.
The odd part below is an arbitrary constant positive integer.

\begin{named}{Lemma 5 (constant odd part).}
\marklemma{5}{lem:fixed-odd}
Let $s$ be a positive integer and set $s' = s/\gcd(s, 3)$.
There is an integer $j_0(s') \ge 11000$ such that for $r \in \{1, 2\}$, every integer $p$, and every integer $j \ge j_0(s')$,
\[
\left| 2^{r/3} - \frac{p}{s\, 2^j} \right| > c_3(s')\, 2^{-\mu_3 j}.
\]
For every $s' \le 16$ one has $j_0(s') = 11000$.
For a larger $s'$ the same integer is $11000 + O(\log s')$.
\end{named}

\begin{proof}
In the argument of Bauer and Bennett, $s \mid 3$ is used once.
It absorbs $s$ into the power of $3$ in the denominator of their linear form $\Lambda$, and gives $|\Lambda| \ge 1/D$.
For a general positive integer $s$, that denominator gains the factor $s' = s/\gcd(s, 3)$, so $|\Lambda| \ge 1/(s' D)$.
Their upper bound and their Proposition~5.1 are unchanged, and so is $\mu_3$.
The factor $1/s'$ divides the constant, which becomes $c_3(s')$.
The choice $r = 2$ in the exponent $N^{-(n+77)/63}$, with $n = 3j+r$, contributes $N^{-79/63}$ after the power of $2^j$ is removed, and it is the smaller of the two choices $r \in \{1, 2\}$.

The term they subtract from the main term is exponentially smaller than $1/D$.
It stays smaller than $1/(2s'D)$ for every $j \ge 11000$ when $s' \le 16$.
For a larger $s'$ the same comparison holds once $j \ge j_0(s')$, with $j_0(s') = 11000 + O(\log s')$.
Both the constant and this integer are read off from their estimates.

Their appeal to the nonvanishing lemma of Beukers (1980/81) uses $s \mid 3$.
For a general $s$, one of the two Padé choices $\delta \in \{0, 1\}$ still gives $\Lambda \ne 0$.
If both choices vanished, $2^{1/3}$ would be rational.
\end{proof}

Per unit of $j$ the restricted irrationality exponent is $\mu_3 \approx 1.61487$, which lies below $\varphi$.

\section{The integer $Q_3(x)$}
\label{sec:q3}

Let $x$ be an odd positive integer.
Let $s'$ be the positive integer $x/\gcd(x, 3)$.
These formulas are the length-$3$ case of Theorem~\ref{thm:bound}, with the factor $1+1/x$ in place of $B = 2 + 1/x$, the constant real number $\mu = \mu_3$, the constant positive real number $c = c_3(s')$, the constant positive real number $c_* = c_3(x)$, and the positive integer $j_0 = J_3(x)$.

Set
\[
A_3 = \frac{x}{2}\, 2^{1/3}
\Biggl( \frac{c_3(s')}{2^{2/3}\,(1+1/x)} \Biggr)^{1/\mu_3}
\]
and
\[
Q_{3,*}(x) = \Biggl( \frac{2/(3A_3) + 3}{c_3(x)\, 2^{1-\mu_3}} \Biggr)^{1/\gamma_3}.
\]
Let $J_3(x) = j_0(s')$.
For each $r \in \{1, 2\}$ and each integer $j$ with $0 \le j < J_3(x)$, and for each odd positive integer $y$ with $y < 3x \cdot 2^{J_3(x)}$, set
\[
\rho = \Biggl| \frac{y}{x\, 2^{j}} \cdot 2^{-r/3} - 1 \Biggr|.
\]
Whenever $\rho > 0$, set the corresponding value equal to $(1+1/x)/\rho$.
The integer $Q_{3,0}(x)$ is the maximum of those finitely many values.
Then
\[
Q_3(x) = \max\bigl(Q_{3,0}(x),\, Q_{3,*}(x)\bigr).
\]

The factor $x/2$ in $A_3$ is the condition $|\rho| < 1/2$.
The power $2^{1/3}$ is the smaller cube root in the lower bound for another element of the cycle, and $2^{2/3}$ is the larger cube root in the error of that approximation.
The numerator $2/(3A_3)+3$ is two other elements, divided by the length, plus the length.

\begin{named}{Corollary 6.}
\marklemma{6}{cor:q3}
Let $x$ be an odd positive integer.
If $x$ lies on a cycle of length $3$, then $q < Q_3(x)$.
\end{named}

\begin{proof}
The statement of Theorem~2 of~\cite{hackman-i} does not apply: the cycle need not visit $1$, and the odd part is $x$ or $3x$ rather than $1$ or $\ell$.
The comparison in its proof does apply, with that odd part.

Consider a cycle of length $3$ through $x$, written $x \to y \to z \to x$.
Lemma~\ref{lem:any-window} gives $2^h = (q+1/x)(q+1/y)(q+1/z)$, so $q = \lfloor 2^{h/3} \rfloor$ and $3 \nmid h$.
Write $2^{h/3} = q + \delta$ with $0 < \delta < 1$, and write $h = 3m + r$ with $r \in \{1, 2\}$.
The fraction $(3xq+1)/(3x \cdot 2^m)$ then satisfies
\[
\Biggl| 2^{r/3} - \frac{3xq+1}{3x \cdot 2^m} \Biggr| = \frac{|\delta - 1/(3x)|}{2^m}.
\]
The denominator $3x \cdot 2^m$ has odd part $3x$, so Lemma~\ref{lem:fixed-odd} applies with constant $c_3(3x/\gcd(3x, 3)) = c_3(x)$.

Lemma~\ref{lem:any-window} gives $e^{-3/q} \le 3\delta/\sigma \le e^{3/q}$ with $\sigma = 1/x + 1/y + 1/z$.
Therefore $|\delta - \sigma/3| \le (\sigma/q)\, e^{3/q}$, and
\[
\bigl|\delta - 1/(3x)\bigr| \le (1/y + 1/z)/3 + (\sigma/q)\, e^{3/q}.
\]
From $qx + 1 = 2^{k_1} y$ and $2^{h/3} = q+\delta$,
\[
\frac{y}{x} = 2^{h/3 - k_1} \cdot \frac{q + 1/x}{q+\delta}.
\]
Since $3 \nmid h$, the exponent $h - 3k_1$ is congruent to $1$ or $2$ modulo $3$.
Write it as $3j + r'$ with $r' \in \{1, 2\}$ and
\[
\frac{y}{x\, 2^{j}} = 2^{r'/3}(1+\rho), \qquad \rho = \frac{1/x - \delta}{q+\delta}.
\]
Then $|\rho| < (1+1/x)/q$.
If $q < 30$, then $q < Q_3(x)$, because Appendix~\ref{app:computations} records $Q_{3,*}(x) > 2^{17906}$.
Assume $q \ge 30$.
Let $\sigma = 1/x + 1/y + 1/z$.
From Lemma~\ref{lem:any-window}, $e^{-3/q} \le 3\delta/\sigma \le e^{3/q}$.
Hence $|3\delta/\sigma - 1| \le e^{3/q} - 1 \le (3/q)\, e^{3/q}$, and $|\delta - \sigma/3| \le (\sigma/q)\, e^{3/q}$.
From $qx + 1 = 2^{k_1} y$ and $qy + 1 = 2^{k_2} z$ there is a real number $\rho_z$ with
\[
\rho_z = \frac{q\bigl(1/x + 1/y - 2\delta\bigr) + 1/(xy) - \delta^2}{(q+\delta)^2}.
\]
Since $1/x$, $1/y$ and $1/z$ lie in $(0, 1]$, one has $|1/x + 1/y - 2/z| \le 2$, so
\[
\bigl|1/x + 1/y - 2\delta\bigr| \le \frac23 + \frac{6\, e^{3/q}}{q}.
\]
Hence $|\rho_z| \le (2/3)/q + \bigl(6e^{3/q} + 2\bigr)/q^2$.
For $q \ge 30$ one has $e^{3/q} < 1.1052$ and $6e^{3/q}+2 < 8.64$, while $8.64/q \le 8.64/30 < 1/3$.
Thus $|\rho_z| < (1+1/x)/q$.
For $q \ge 30$ both $|\rho|$ and $|\rho_z|$ are less than $1/2$.
If $|\rho| \ge 1/2$, then $q < 2(1+1/x)$.
If $|\rho| < 1/2$ and $j \ge J_3(x)$, Lemma~\ref{lem:fixed-odd} and $|2^{r'/3} - y/(x\, 2^j)| = 2^{r'/3}|\rho| \le 2^{2/3}|\rho|$ give
\[
2^{j} > \Biggl( \frac{c_3(s')\, q}{2^{2/3}\,(1+1/x)} \Biggr)^{1/\mu_3},
\]
and therefore $y > A_3\, q^{1/\mu_3}$.
The same lower bound holds for $z$.
Then $\sigma \le 1/x + 2 A_3^{-1} q^{-1/\mu_3}$, so
\[
(\sigma/q)\, e^{3/q} \le \frac{e^{3/q}}{xq} + 2 A_3^{-1} e^{3/q}\, q^{-1-1/\mu_3}.
\]
Here $N < 2^{34}$ and $79/63 < 4/3$, so $N^{79/63} < 2^{46}$ and $c_3(s') > 2^{-51}/x$.
Also $N > 2^{33}$, so $\mu_3 > 3/2$, hence $1/\mu_3 < 2/3$, and $2^{2/3}(1+1/x) \le 4$.
Hence $A_3 > 2^{-37}\, x^{1/3}$ and $4A_3^{-1} < 2^{39}$.
Appendix~\ref{app:computations} records $Q_{3,*}(x) > 2^{17906}$, so $q > 4A_3^{-1}$ whenever $q > Q_{3,*}(x)$.
For such $q$ one has $e^{3/q} < 2$ and $q^{-1} \le q^{-1/\mu_3}$, because $\mu_3 > 1$.
The first summand is then at most $2\, q^{-1/\mu_3}$, and the second is at most $q^{-1/\mu_3}$.
Therefore
\[
\bigl|\delta - 1/(3x)\bigr| \le \bigl(2/(3A_3) + 3\bigr)\, q^{-1/\mu_3}.
\]
If $m \ge J_3(x)$ and $\mu_3 > 1$, the lower bound from Lemma~\ref{lem:fixed-odd}, together with $2^m < q+1$, gives
\[
\bigl|\delta - 1/(3x)\bigr| > c_3(x)\, 2^{1-\mu_3}\, q^{1-\mu_3}.
\]
These two bounds on $|\delta - 1/(3x)|$ conflict once $q > Q_{3,*}(x)$.
If $s' \le 16$, then $J_3(x) = 11000$, so $q > Q_{3,*}(x)$ forces $q > 2^{11001}$ and $m \ge 11000$.
If $s' > 16$ and $q \le 2^{J_3(x)+1}$, that power of $2$ is an upper bound.
It does not change $\max\bigl(Q_{3,0}(x), Q_{3,*}(x)\bigr)$ whenever $J_3(x) \le 17905$, and $J_3(x) = 11000$ for every $s' \le 16$.

If instead $0 \le j < J_3(x)$, the value $(1+1/x)/\rho$ is one of the numbers whose maximum is $Q_{3,0}(x)$, so $q \le Q_{3,0}(x)$.
Thus $q < Q_3(x)$.
\end{proof}

\section{The integer $Q_2(x)$}
\label{sec:q2}

These formulas are the length-$2$ case of Theorem~\ref{thm:bound}, with the factor $1+1/x$ in place of $B = 1$, the constant real number $\mu = 1.465$, the constant positive real numbers $c = c_* = 1$, and the positive integer $j_0 = j_0(x)$.

Let $x$ be an odd positive integer.
Let $\mu_2 = 1.465$ and $\gamma_2 = 1 - \mu_2 + 1/\mu_2$.
The value $\mu_2 = 1.465$ is the exponent printed to three decimals in Corollary~1.4 of Bauer and Bennett~\cite{bauer-bennett-2002}, for the square root of $2$.
For each constant positive integer $s$ that corollary gives an integer $j_0(s)$ such that for every integer $p$ and every integer $j \ge j_0(s)$,
\[
\Bigl| \sqrt{2} - \frac{p}{s\, 2^{j}} \Bigr| > 2^{-\mu_2 j}.
\]
The constant in front is $1$.
The integer $j_0(x)$ is that exponent for $s = x$.
Set
\[
A_2 = \frac{x}{2}\, 2^{1/2}
\Biggl( \frac{1}{2^{1/2}\,(1+1/x)} \Biggr)^{1/\mu_2}
\]
and
\[
Q_{2,*}(x) = \Biggl( \frac{1/(2A_2) + 2}{2^{1-\mu_2}} \Biggr)^{1/\gamma_2}.
\]
For each integer $j$ with $0 \le j < j_0(x)$, and for each odd positive integer $y$ with $y < 3x \cdot 2^{j_0(x)}$, set
\[
\rho = \Biggl| \frac{y}{x\, 2^{j}} \cdot 2^{-1/2} - 1 \Biggr|.
\]
Whenever $\rho > 0$, set the corresponding value equal to $(1+1/x)/\rho$.
The integer $Q_{2,0}(x)$ is the maximum of those finitely many values.
The length-$2$ bound is this integer:
\[
Q_2(x) = Q_{2,0}(x).
\]
The quantities $2^{j_0(x)+1}$ and $Q_{2,*}(x)$ are at most $Q_{2,0}(x)$ for every odd positive integer $x$.
The proof below records that comparison.

\begin{named}{Corollary 7.}
\marklemma{7}{cor:q2}
Let $x$ be an odd positive integer.
If $x$ lies on a cycle of length $2$, then $q < Q_2(x)$.
\end{named}

\begin{proof}
Consider a cycle of length $2$ through $x$, written $x \to y \to x$.
Here $h$ is odd.
Write $2^{h/2} = q + \delta$ and $h = 2m+1$.
The argument for $Q_3(x)$ in Section~\ref{sec:q3}, with one other element and with Corollary~1.4 of Bauer and Bennett (2002)~\cite{bauer-bennett-2002} in place of Lemma~\ref{lem:fixed-odd}, produces $A_2$ and $Q_{2,*}(x)$.
Both the product around the cycle and the approximation for $y$ have odd part $x$, and the constant in front of $2^{-\mu_2 j}$ is $1$.
The numerator $1/(2A_2)+2$ is that one other element, divided by the length, plus the length.

The exponent $j = 0$ fails the inequality of Corollary~1.4 for every $x \ge 1$, so $j_0(x) \ge 1$.
The distance from $x\, 2^{j_0(x)-1}\sqrt{2}$ to an odd integer is at most $1$, and the resulting value of $(1+1/x)/\rho$ is at least $(x+1)\sqrt{2}\, 2^{j_0(x)-1}$.
For $x \ge 3$ this is at least $2^{j_0(x)+1}$, because $(x+1)\sqrt{2} \ge 4\sqrt{2} > 4$.
For $x = 1$ one has $j_0(1) = 9$, and the value at $j = 7$ has logarithm $14.192568 > 10$.
Thus $2^{j_0(x)+1} \le Q_{2,0}(x)$.

The quantity $A_2$ increases with $x$, so $Q_{2,*}(x)$ decreases, and $Q_{2,*}(1) < 2^{11}$.
Every exponent $j$ with $x \cdot 2^{-0.465 j} \ge 1/2$ fails the inequality, and the largest such $j$ contributes at least $(1+1/x)\, x \sqrt{2}\, 2^{j}$ even when the distance to an odd integer is $1$.
For every odd $x \ge 5$ that lower bound already exceeds $Q_{2,*}(x)$.
For $x = 1$ and $x = 3$ the values $Q_{2,0}(1) > 2^{14}$ and $Q_{2,0}(3) > 2^{17}$ do the same.
Thus $Q_{2,*}(x) \le Q_{2,0}(x)$, and $Q_2(x) = Q_{2,0}(x)$.
For $\ell = 2$, Theorem~\ref{thm:bound} has $B = 1$ and $Q_A = 1$.
The factor $1+1/x$ is at least $B$, so each value $B/\rho$ is at most the corresponding value $(1+1/x)/\rho$.
At $q = 8$ one has $q \ge 4\cdot 2$ and $e^{2/8} < 8^{1 - 1/1.465}$, so $Q_e \le 8$.
The exponent $j = 2$ satisfies $x \cdot 2^{-0.465 \cdot 2} > 1/2$ for every $x \ge 1$, so it fails the inequality, and the value $(1+1/x)/\rho$ is at least $4(x+1)\sqrt{2} \ge 8\sqrt{2} > 8$.
Thus $2^{j_0(x)}$ and $Q_e$ are at most $Q_{2,0}(x)$.
The remaining values in $Q_0$ are $\sqrt{2}\, x / d$, where $d$ is the distance from $x\sqrt{2}$ to the nearest integer.
If $d \ge 1$, that value is at most $\sqrt{2}\, x$, which is less than $8\sqrt{2} \le Q_{2,0}(x)$ for $x = 1$ and less than $4(x+1)\sqrt{2}$ for $x \ge 3$.
If $d < 1$ and $p$ is the nearest integer, then $|p^2 - 2x^2| \ge 1$ and $p + x\sqrt{2} < 2x\sqrt{2} + 1 < 4x$, so $d > 1/(4x)$ and the value is less than $4\sqrt{2}\, x^2$.
For $x = 1$ one has $4\sqrt{2} < 2^{14} < Q_{2,0}(1)$.
For $x \ge 3$, the largest integer $j$ with $x \cdot 2^{-0.465 j} \ge 1/2$ satisfies $2^j \ge (2x)^{1/0.465}/2$, and
\[
(x+1)\sqrt{2}\cdot 2^{j} \ge \frac{\sqrt{2}\, x\, (2x)^{1/0.465}}{2}.
\]
The quotient of the right side by $4\sqrt{2}\, x^2$ is $2^{1/0.465}\, x^{1/0.465 - 1}/8 > 2^{2.15}\, x / 8$.
Since $2^{2.15} > 4.4$, the quotient is greater than $0.55x \ge 0.55 \cdot 3 > 1$.
Thus the maximum in Theorem~\ref{thm:bound} is at most $Q_2(x)$, and a cycle of length $2$ through $x$ satisfies $q < Q_2(x)$.

For $x = 1$ and length $2$, Proposition~A.1 of~\cite{hackman-i} gives the only such $q$, namely $q = 5$.
That conclusion does not change the formula for $Q_2(1)$.
\end{proof}

\section{Cycles of length $2$}
\label{sec:length-2}

The aim of the research in this section was to prove that the only cycles of length $2$ are $1 \to 3$, $35 \to 99$, and $27 \to 611$.
That is more ambitious than Theorem~1 of the companion paper, which only rules out a cycle of length $3$ through $1$.
Here the starting value may be any odd positive integer.
A cycle of length $2$ was selected because it's the smallest non-trivial length, making it look like the easiest problem of this kind.

For a constant starting value, Theorem~\ref{thm:bound} bounds $q$, and a computation can finish the search.
The bound grows with the starting value, so the same search does not cover every starting value at once.
Proposition~\ref{prop:exponent} shows that a restricted irrationality exponent of $1.3247$, for every odd positive integer, gives a bound on $q$ for every starting value and every length at least $2$.
Bauer and Bennett (2002)~\cite{bauer-bennett-2002} give the exponent $1.465$ for $\sqrt{2}$.

The aim was not achieved, but this section records the properties of a cycle of length $2$, and large bounds that follow from those properties by a modest computation.

Throughout this section, $x \to y \to x$ is a cycle of length $2$, and $x$ is its smallest element.
Here $qx + 1 = 2^{k_1} y$ and $qy + 1 = 2^{k_2} x$, and $h = k_1 + k_2$.
Lemma~\ref{lem:any-window} gives $q = \lfloor 2^{h/2} \rfloor$ and $h$ odd.
Let $D = 2^h - q^2$.

\begin{named}{Lemma 8 (the length-$2$ equation).}
\marklemma{8}{lem:length-2}
Then $k_1 < k_2$.
Let $t = k_2 - k_1$.
Then $t$ is positive and odd, $h = 2k_1 + t$, and
\[
x = \frac{q + 2^{k_1}}{D}, \qquad y = \frac{q + 2^{k_2}}{D}.
\]
The integer $D$ divides $2^t - 1$, and $D \equiv 7 \pmod{8}$.
Every prime divisor of $D$ is congruent to $1$ or $7$ modulo $8$.
Let $u = (2^t - 1)/D$.
Then $u$ is a positive odd integer,
\[
y = x + u\, 2^{k_1}, \qquad y^2 - 2^t x^2 = u,
\]
$2^{k_1} \le q - 1$, $t > \log_2(q/x)$, and $y^2 > qx + 1$.
\end{named}

\begin{proof}
Since $h$ is odd, $k_1 \ne k_2$.
From $qy + 1 = 2^{k_2} x$ and $y = (qx + 1)/2^{k_1}$,
\[
2^h x = q(qx + 1) + 2^{k_1},
\]
so $x(2^h - q^2) = q + 2^{k_1}$.
Lemma~\ref{lem:any-window} gives $q^2 < 2^h$, so $D > 0$ and $x = (q + 2^{k_1})/D$.
The same elimination with the roles of $x$ and $y$ exchanged gives $y = (q + 2^{k_2})/D$.
Since $x < y$ and $D > 0$, one has $k_1 < k_2$.
Thus $t > 0$.
From $h = 2k_1 + t$ and $h$ odd, $t$ is odd.

Now $D$ divides $q + 2^{k_1}$ and $q + 2^{k_2}$, so $D$ divides
\[
(q + 2^{k_2}) - (q + 2^{k_1}) = 2^{k_1}(2^t - 1).
\]
Here $q$ is odd and $h \ge 3$, so $2^h$ is divisible by $8$ and $q^2 \equiv 1 \pmod{8}$.
Thus $D \equiv 7 \pmod{8}$, and in particular $D$ is odd.
Hence $D$ divides $2^t - 1$.
Let $p$ be a prime divisor of $D$.
Then $q^2 = 2^h - D \equiv 2^h \pmod{p}$.
If $p$ divides $q$, then $p$ divides $2^h$.
But $D$ is odd, so $p \ne 2$, and $p$ does not divide $2^h$.
Thus $2^h$ is a nonzero square modulo $p$.
Since $h$ is odd, $2$ is a quadratic residue modulo $p$, and therefore $p \equiv 1 \pmod{8}$ or $p \equiv 7 \pmod{8}$.

The quotient $u = (2^t - 1)/D$ is a positive integer.
Both $2^t - 1$ and $D$ are odd, so $u$ is odd.
Then
\[
y - x = \frac{2^{k_2} - 2^{k_1}}{D} = u\, 2^{k_1},
\]
so $y = x + u\, 2^{k_1}$.
Substitute $q = Dx - 2^{k_1}$ into $qx + 1 = 2^{k_1} y$:
\[
D x^2 - 2^{k_1 + 1} x + 1 = u\, 4^{k_1}.
\]
Therefore
\begin{align*}
y^2 - 2^t x^2
&= \bigl(x + u\, 2^{k_1}\bigr)^2 - 2^t x^2 \\
&= u\, 2^{k_1 + 1} x + u^2 4^{k_1} - x^2(2^t - 1) \\
&= u\bigl(2^{k_1 + 1} x + u\, 4^{k_1} - D x^2\bigr) \\
&= u.
\end{align*}

Since $x$ and $y$ are distinct and odd, $y \ge x + 2$.
Then
\[
2^{k_1} = \frac{qx + 1}{y} \le \frac{qx + 1}{x + 2}.
\]
The inequality $(qx + 1)/(x + 2) < q$ is $1 < 2q$.
Thus $2^{k_1} < q$.
Both are integers, so $2^{k_1} \le q - 1$.

From $y - x \ge 2^{k_1}$ and $2^{k_1} = (qx + 1)/y$,
\[
2^{k_1}\bigl(x + 2^{k_1}\bigr) \le qx + 1,
\]
so $2^{2k_1} \le x(q - 2^{k_1}) + 1$.
For $k_1 \ge 1$ the right side is at most $x(q - 2) + 1$, and $x(q - 2) + 1 < xq$ because $1 < 2x$.
Thus $2^{2k_1} < xq$.
Lemma~\ref{lem:any-window} gives $2^{h/2} = q + \delta$ with $\delta > 0$, so $h > 2\log_2 q$.
Therefore
\[
t = h - 2k_1 > 2\log_2 q - \log_2(xq) = \log_2(q/x).
\]

Finally $y(y - x) \ge y \cdot 2^{k_1} = qx + 1$.
Since $x \ge 1$, one has $y > y - x$, and therefore $y^2 > qx + 1$.
\end{proof}

\begin{named}{Lemma 9 (the case $D = 7$).}
\marklemma{9}{lem:gap-7}
If $D = 7$, the cycle is $1 \to 3$ with $q = 5$, or $35 \to 99$ with $q = 181$, or $27 \to 611$ with $q = 181$.
\end{named}

\begin{proof}
Here $q^2 + 7 = 2^h$.
By Theorem~1.8 of Bauer and Bennett (2002), the positive solutions $(q, h)$ are $(1, 3)$, $(3, 4)$, $(5, 5)$, $(11, 7)$ and $(181, 15)$.
Lemma~\ref{lem:length-2} requires $q \ge 3$ and $h$ odd, which leaves $(5, 5)$, $(11, 7)$ and $(181, 15)$.
The order of $2$ modulo $7$ is $3$, so $3$ divides $t$, and Lemma~\ref{lem:length-2} requires $t$ odd and $1 \le t \le h - 2$.

For $(q, h) = (5, 5)$ the only such $t$ is $3$, so $k_1 = 1$ and $k_2 = 4$.
Then $(5 + 2^1)/7 = 1$ and $(5 + 2^4)/7 = 3$, which is $1 \to 3$.

For $(q, h) = (11, 7)$ one has $-11 \equiv 3 \pmod{7}$.
The powers of $2$ modulo $7$ are $2$, $4$ and $1$, so $q + 2^{k_1}$ is never divisible by $7$.

For $(q, h) = (181, 15)$ the admissible values are $t = 3$ and $t = 9$.
If $t = 3$, then $k_1 = 6$ and $k_2 = 9$, and
\[
\frac{181 + 2^6}{7} = 35, \qquad \frac{181 + 2^9}{7} = 99,
\]
which is $35 \to 99$.
If $t = 9$, then $k_1 = 3$ and $k_2 = 12$, and
\[
\frac{181 + 2^3}{7} = 27, \qquad \frac{181 + 2^{12}}{7} = 611,
\]
which is $27 \to 611$.
\end{proof}

\begin{named}{Lemma 10 (the case $u = 1$).}
\marklemma{10}{lem:u-one}
If $u = 1$, the cycle is $1 \to 3$ with $q = 5$, or $35 \to 99$ with $q = 181$.
\end{named}

\begin{proof}
Here $D = 2^t - 1$, so $q^2 = 2^{2k_1 + t} - 2^t + 1$ and
\[
(q - 1)(q + 1) = 2^t\bigl(2^{2k_1} - 1\bigr).
\]
The factor $2^{2k_1} - 1$ is odd, and $\gcd(q - 1, q + 1) = 2$, so
\[
v_2(q - 1) + v_2(q + 1) = t
\]
with the smaller valuation equal to $1$.
Lemma~\ref{lem:length-2} gives $t$ odd and $D \equiv 7 \pmod{8}$, so $t \ge 3$.
Let $c = t - 2$, which is odd and at least $1$.

If $v_2(q + 1) = 1$ and $v_2(q - 1) = t - 1$, write $q + 1 = 2r$ and $q - 1 = 2^{t - 1} s$ with $r$ and $s$ odd.
Then $r = 2^{t - 2} s + 1$ and $s(2^{t - 2} s + 1) = 2^{2k_1} - 1$, so
\[
2^c s^2 + s + 1 = 2^{2k_1}.
\]
If $v_2(q - 1) = 1$ and $v_2(q + 1) = t - 1$, write $q - 1 = 2s$ and $q + 1 = 2^{t - 1} r$ with $r$ and $s$ odd.
Then $s = 2^{t - 2} r - 1$ and $r(2^{t - 2} r - 1) = 2^{2k_1} - 1$, so
\[
2^c r^2 - r + 1 = 2^{2k_1}.
\]
In either case, for a positive odd integer $X$,
\[
\bigl(2^{c + 1} X \pm 1\bigr)^2 + \bigl(2^{c + 2} - 1\bigr) = 2^{2k_1 + c + 2},
\]
the plus sign for $X = s$ and the minus sign for $X = r$.

Suppose $c \ge 3$.
Then $2^{c + 2} - 1 = 2^k - 1$ with $k = c + 2 \ge 5$.
Theorem~1.8 of Bauer and Bennett (2002) gives exactly two positive solutions of $Y^2 + (2^k - 1) = 2^n$, namely $Y = 1$ with $n = k$, and $Y = 2^{k - 1} - 1$ with $n = 2k - 2$.
The first gives $2k_1 + c + 2 = c + 2$, so $k_1 = 0$.
The second gives $2k_1 + c + 2 = 2c + 2$, so $2k_1 = c$.
But $c$ is odd, so neither solution has an integer $k_1 \ge 1$.

Now $c = 1$, so $t = 3$ and $D = 7$, and the equation is $Y^2 + 7 = 2^{2k_1 + 3}$ with $Y = 4X \pm 1$.
The same theorem gives $Y \in \{1, 3, 5, 11, 181\}$ and $2k_1 + 3 \in \{3, 4, 5, 7, 15\}$.
The exponent $2k_1 + 3$ is odd and at least $5$, which leaves $5$, $7$ and $15$.
If the exponent is $5$, then $Y = 5$ and $k_1 = 1$.
The equation $4X - 1 = 5$ has no integer solution, while $4X + 1 = 5$ gives $X = 1$.
This is the plus sign, $q - 1 = 2^{2} \cdot 1$, and $q = 5$, the cycle $1 \to 3$.
If the exponent is $7$, then $Y = 11$ and $k_1 = 2$.
The equation $4X + 1 = 11$ has no integer solution, while $4X - 1 = 11$ gives $X = 3$.
This is the minus sign and $q = 11$.
Then $(11 + 2^2)/7 = 15/7$, which is not an integer.
If the exponent is $15$, then $Y = 181$ and $k_1 = 6$.
The equation $4X - 1 = 181$ has no integer solution, while $4X + 1 = 181$ gives $X = 45$.
This is the plus sign, $q - 1 = 2^{2} \cdot 45$, and $q = 181$, the cycle $35 \to 99$.
\end{proof}

\begin{named}{Proposition 11.}
\marklemma{11}{prop:other-two}
Any cycle of length $2$ other than $1 \to 3$, $35 \to 99$ and $27 \to 611$ has
\[
q \ge 2^{160{,}000}, \qquad D \ge 15, \qquad u \ge 3,
\]
every odd element at least $37$, and, writing $h = 2m + 1$,
\[
m \ge 160{,}000, \qquad D > 2^{0.535 m + 1}, \qquad t \ge 85{,}603.
\]
In particular $k_1 \le \lfloor \log_2 q \rfloor - 42{,}801$ and $k_2 \ge \lfloor \log_2 q \rfloor + 42{,}802$.
Also $t > \log_2(3q/x)$.
\end{named}

\begin{proof}
Appendix~\ref{app:computations} searches every cycle of length $2$ with $h \le 320{,}000$ and finds only the three cycles named above.
That search excludes every cycle of length $2$ with $q < 2^{160{,}000}$.
Thus any other cycle has $q \ge 2^{160{,}000}$.
Since $2^{h/2} > q$, one has $h > 2 \cdot 160{,}000$, so $h \ge 320{,}001$.
With $h = 2m + 1$ this is $m \ge 160{,}000$.

Lemma~\ref{lem:gap-7} gives $D \ne 7$.
Lemma~\ref{lem:length-2} gives $D \equiv 7 \pmod{8}$ and $D > 0$, so $D \ge 15$.
Lemma~\ref{lem:u-one} gives $u \ne 1$.
Lemma~\ref{lem:length-2} gives that $u$ is a positive odd integer, so $u \ge 3$.

Let $z$ be an odd element of the cycle with $z \le 35$.
Theorem~\ref{thm:bound} gives $q < Q_2(z)$.
Appendix~\ref{app:computations} records $Q_2(z) \le 2^{33}$ for every odd $z \le 35$.
Then $q < 2^{33} < 2^{160{,}000}$, so the cycle is one of the three found by the search.
Each of those has its odd elements in $\{1, 3, 27, 35, 99, 611\}$.
An element $z \le 35$ on one of them is $1$, $3$, $27$ or $35$.
Therefore every odd element of any other cycle is at least $37$.

Corollary~1.4 of Bauer and Bennett (2002), applied with odd part $1$, gives
\[
\Bigl| \sqrt{2} - \frac{q}{2^m} \Bigr| > 2^{-1.465 m}
\]
for every $m \ge j_0(1)$.
Section~\ref{sec:q2} records $j_0(1) = 9$, and here $m \ge 160{,}000$.
Writing $2^{h/2} = q + \delta$ as in Lemma~\ref{lem:any-window},
\[
\delta > 2^{-0.465 m}.
\]
Then $D = \delta(2q + \delta) > 2q\delta$, so
\[
\log_2 D > 1 + \log_2 q - 0.465 m.
\]
Since $q > 2^m \sqrt{2} - 1$ and $m \ge 160{,}000$,
\[
\log_2 q > m + \tfrac12 - 2^{-159{,}999}.
\]
Hence $\log_2 D > 0.535 m + 1.5 - 2^{-159{,}999} > 0.535 m + 1$, and $D > 2^{0.535 m + 1}$.
At $m = 160{,}000$ this is $D > 2^{85{,}601}$.
For $m > 160{,}000$ the same lower bound is larger.
Lemma~\ref{lem:length-2} gives $D \le 2^t - 1$, so $t > 85{,}601$.
Since $t$ is odd, $t \ge 85{,}603$.

Since $t$ is odd, $(t - 1)/2$ is an integer, and $k_1 = m - (t - 1)/2$.
Thus $k_1 \le m - 42{,}801$.
Since $m \ge 160{,}000$ and $\sqrt{2} - 1 > 2^{-m}$, one has $2^m \sqrt{2} > 2^m + 1$, so $q \ge 2^m + 1$.
Also $\sqrt{2} < 2$, so $q < 2^m \sqrt{2} < 2^{m + 1}$.
Thus $\lfloor \log_2 q \rfloor = m$ and $k_1 \le \lfloor \log_2 q \rfloor - 42{,}801$.
Also $k_2 = m + (t + 1)/2 \ge m + 42{,}802 = \lfloor \log_2 q \rfloor + 42{,}802$.

For the last claim, $u \ge 3$ gives $y - x \ge 3 \cdot 2^{k_1}$.
Combined with $2^{k_1} = (qx + 1)/y$,
\[
3 \cdot 4^{k_1} \le x\bigl(q - 2^{k_1}\bigr) + 1 < xq,
\]
so $2^{2k_1} < xq/3$.
With $h > 2\log_2 q$,
\[
t = h - 2k_1 > 2\log_2 q - \log_2(xq/3) = \log_2(3q/x).
\]
\end{proof}

\begin{named}{Lemma 12.}
\marklemma{12}{lem:onset-three}
Any cycle of length $2$ other than $1 \to 3$, $35 \to 99$ and $27 \to 611$ has every odd element greater than $2^{2559}$.
\end{named}

\begin{proof}
Suppose the cycle is other than those three, and suppose $x \le 2^{2559}$.

Theorem~1.5 of Bauer and Bennett (2002)~\cite{bauer-bennett-2002} applies to the solution $181^2 + 7 = 2^{15}$, with base $2$, with $a = 1$, with the power of $2$ equal to $15$, with constant term $7$ and with $\alpha = 3$, taking the odd part to be $x$.
In that theorem $\log$ is the natural logarithm.
Let $\xi = 7/2^{15}$.
Formula~(1.2) of that paper gives $r(3, \xi) = 0.250030044752$, and formula~(1.3) gives $F(r(3, \xi), 3, \xi) = 9.479962374117$.
Table~(1.17) gives $c_1(3) = 2.458$ and $\log d_1(3) = -17.335$.
The remaining parameters in that theorem are $\Delta_0 = 1$, $\kappa(7) = 4$, $\Omega = 2^{17}$, $m_1 = 17$, $\Omega_1 = 2^{17}$ and $\Delta_1 = 1$.
Formulas~(1.13) and~(1.16) do not depend on $x$.
They give
\[
\chi_1 = 4.968969066213, \qquad \chi_3 = 8.684694 \times 10^{15}
\]
and
\[
\lambda_1 = \frac{\log \chi_3}{2 \cdot 17 \cdot \log 2} = 1.5572763652.
\]
Let $\lambda_2 = 1.617$ be a constant real number.
Then $\lambda_2 < \varphi$.

Let $c_2(x)$ be the constant in Theorem~1.5 at odd part $x$, and let $m_2(x)$ be the minimum in (1.15) at the prime $2$.
Define the real numbers
\[
k(x) = \frac{\log\bigl(1/c_2(x)\bigr)}{(\lambda_2 - \lambda_1)\log 2}, \qquad k_0(x) = \max\bigl(m_2(x) - \tfrac12,\, k(x)\bigr).
\]
Theorem~1.5 gives, for every integer $p$ and every integer $j \ge m_2(x) - 1/2$,
\[
\Bigl| \sqrt{2} - \frac{p}{x\, 2^{j}} \Bigr| > c_2(x)\, 2^{-\lambda_1 j}.
\]
For every integer $j \ge k(x)$ one has $c_2(x)\, 2^{-\lambda_1 j} \ge 2^{-\lambda_2 j}$.
Therefore, for every integer $p$ and every integer $j \ge k_0(x)$,
\[
\Bigl| \sqrt{2} - \frac{p}{x\, 2^{j}} \Bigr| > 2^{-\lambda_2 j}.
\]
Formula~(1.14), at $\alpha = 3$ and at odd part $x$, increases with $x$, and $c_2(x)$ decreases as $x$ increases.
Both $m_2(x)$ and $k(x)$ increase with $x$, and so does $k_0(x)$.
At $x = 1$ one has $k_0(1) = 2840.688\ldots$, so $|\sqrt{2} - p/2^{j}| > 2^{-\lambda_2 j}$ for every integer $j \ge 2841$.
At $x = 2^{2559}$ one has $k_0(2^{2559}) = 77426.543\ldots$.
Since $x \le 2^{2559}$, one has $k_0(x) \le k_0(2^{2559})$.

Let $\gamma = 1 - \lambda_2 + 1/\lambda_2$ and
\[
A(x) = \frac{x}{2}\, 2^{1/2}
\Biggl( \frac{1}{2^{1/2}\,(1+1/x)} \Biggr)^{1/\lambda_2},
\qquad
Q_*(x) = \Biggl( \frac{1/(2A(x)) + 2}{2^{1-\lambda_2}} \Biggr)^{1/\gamma}.
\]
Here $A(x)$ increases with $x$, so $Q_*(x)$ decreases with $x$, and $\log_2 Q_*(1) < 1651$.

Lemma~\ref{lem:any-window} gives $2^{h/2} = q + \delta$ with $0 < \delta < 1$.
Write $h = 2m + 1$ with $m$ a nonnegative integer, and let $j = m - k_1$.
Lemma~\ref{lem:length-2} gives $k_1 = m - (t-1)/2$, so $j = (t-1)/2$ is a nonnegative integer and $h/2 - k_1 = j + 1/2$.
From $qx + 1 = 2^{k_1} y$,
\[
\frac{y}{x\, 2^{j}} = 2^{1/2}(1 + \rho), \qquad \rho = \frac{1/x - \delta}{q + \delta}.
\]
Then $|\rho| < (1 + 1/x)/q$, and therefore $q < (1 + 1/x)/|\rho|$.

If $j \ge k_0(x)$, the argument in the proof of Theorem~\ref{thm:bound} for length $2$ applies with $\lambda_2$ in place of $1.465$ and with $k_0(x)$ in place of $j_0(x)$.
The constant in front of $2^{-\lambda_2 j}$ is $1$, and $1/(2A(x)) + 2$ is the one other element, divided by the length, plus the length.
That argument gives $q < Q_*(x)$, so $q < 2^{1651}$.

If $j < k_0(x)$, then $j \le 77426$.
If $|y - x\, 2^{j}\, 2^{1/2}| > 1$, then $|\rho| > 1/(x\, 2^{j}\, 2^{1/2})$, so
\[
q < (1 + 1/x)\, x\, 2^{j}\, 2^{1/2} < 3\, x\, 2^{j} \le 3 \cdot 2^{2559 + 77426} < 2^{79987}.
\]
If $|y - x\, 2^{j}\, 2^{1/2}| \le 1$, then $y + x\, 2^{j}\, 2^{1/2} \le 2x\, 2^{j}\, 2^{1/2} + 1 < 4x\, 2^{j}$.
The integers $y$ and $x\, 2^{j}$ satisfy $y^2 \ne 2x^2\, 4^{j}$, since $2^{1/2}$ is irrational.
\[
\Bigl| 2^{1/2} - \frac{y}{x\, 2^{j}} \Bigr| = \frac{|2x^2\, 4^{j} - y^2|}{x\, 2^{j}\, |y + x\, 2^{j}\, 2^{1/2}|} > \frac{1}{4x^2\, 4^{j}}.
\]
Hence $|\rho| > 1/(4 \cdot 2^{1/2}\, x^2\, 4^{j})$ and
\[
q < (1 + 1/x) \cdot 4 \cdot 2^{1/2}\, x^2\, 4^{j} \le 8 \cdot 2^{1/2}\, x^2\, 4^{j} < 12\, x^2\, 4^{j},
\]
so $\log_2 q < 2\log_2 x + 2j + 4 < 2\log_2 x + 2k_0(x) + 4$.
The quantity $2\log_2 x + 2k_0(x) + 4$ increases with $x$, and at $x = 2^{2559}$ it is less than $160{,}000$.
In both subcases $q < 2^{160{,}000}$.

Appendix~\ref{app:computations} excludes every cycle of length $2$ with $q < 2^{160{,}000}$ except $1 \to 3$, $35 \to 99$ and $27 \to 611$.
Each of those has its odd elements in $\{1, 3, 27, 35, 99, 611\}$.
The supposition $x \le 2^{2559}$ would make the cycle one of those three.
Therefore $x > 2^{2559}$.
Every odd element is at least $x$, so every odd element is greater than $2^{2559}$.
\end{proof}

The case that remains is $u \ge 3$.
Proposition~\ref{prop:other-two} puts $D \ge 15$, with $D$ dividing $2^t - 1$ and $(q + 2^{(h - t)/2})/D$ an odd integer.
Lemma~\ref{lem:onset-three} gives that this odd integer is greater than $2^{2559}$.
Proposition~\ref{prop:scan-x} gives the lower bounds on $q$ and on $x$.
Lemma~\ref{lem:linear-two} turns the lower bound on $q$ into the remaining bounds.

\begin{named}{Proposition 13.}
\marklemma{13}{prop:scan-x}
Let the cycle be other than $1 \to 3$, $35 \to 99$ and $27 \to 611$.
Then $q \ge 2^{1{,}071{,}999{,}971}$ and $x > 2^{17{,}631{,}663}$.
\end{named}

\begin{proof}
Proposition~\ref{prop:other-two} gives $m \ge 160{,}000$.
Appendix~\ref{app:computations} records the first $1{,}072{,}000{,}000$ binary places of $\sqrt{2}$.
For every integer $m$ with $160{,}000 \le m \le 1{,}071{,}999{,}970$, places $m+1$ through $m+30$ of that expansion are not all $0$.
Suppose $m \le 1{,}071{,}999{,}970$.
The fractional part $\delta$ in $2^{h/2} = q + \delta$ is read from place $m+1$, so $\delta \ge 2^{-30}$.
Lemma~\ref{lem:length-2} gives $x = (q + 2^{k_1})/D$ and $2^{k_1} \le q - 1$, and $D = \delta(2q + \delta)$, so
\[
\delta = \frac{D}{2q + \delta} < \frac{q + 2^{k_1}}{2qx} \le \frac{2q - 1}{2qx} < \frac{1}{x}.
\]
Thus $2^{-30} \le \delta < 1/x$, and $x < 2^{30}$.
Lemma~\ref{lem:onset-three} gives $x > 2^{2559}$, so $m \ge 1{,}071{,}999{,}971$.
Since $m \ge 160{,}000$ and $\sqrt{2} - 1 > 2^{-m}$, one has $2^m \sqrt{2} > 2^m + 1$, so $q \ge 2^m + 1$.
Therefore $q \ge 2^{1{,}071{,}999{,}971} + 1$.

Suppose $x \le 2^{17{,}631{,}663}$.
In the notation of the proof of Lemma~\ref{lem:onset-three}, $\lambda_2 - \lambda_1 = 0.0597236348$ and $k_0(x) = \max\bigl(m_2(x) - 1/2,\, k(x)\bigr)$.
Formula~(1.14) of Bauer and Bennett (2002) is linear in the odd part, and the denominator of $c_2$ contains one factor of the odd part.
The other parameters in the proof of Lemma~\ref{lem:onset-three} do not depend on $x$.
Thus $c_2(x) = c_2(1)/x$ and $\chi_2(x) = x\, \chi_2(1)$, and
\[
k(x) = k(1) + \frac{\log_2 x}{\lambda_2 - \lambda_1}.
\]
Equation~(1.15) of Bauer and Bennett (2002), at the single prime $2$, with $\mathrm{ord}_2(\Omega_1) = 17$ and $\mathrm{ord}_2(a) = 0$, is
\[
m_2(x) = \frac{68 \log\bigl(x\, \chi_2(1)\bigr)}{\log \chi_1} + 15.
\]
Here $\log$ is the natural logarithm, $\log \chi_1 = \log 4.968969066213$, and the same parameters give $\log\bigl(\chi_2(1)\bigr) = 51.353232\ldots$.
The proof of Lemma~\ref{lem:onset-three} records $k(1) = 2840.688\ldots$.
At $x = 2^{17{,}631{,}663}$ one has $k(x) < 295{,}223{,}704$ and $m_2(x) - 1/2 < 518{,}368{,}291$, so $k_0(x) < 518{,}368{,}291$.
Since $k_0$ increases with the odd part, $k_0(x) < 518{,}368{,}291$ for every odd part at most $2^{17{,}631{,}663}$.
Let $j = m - k_1$.
Lemma~\ref{lem:length-2} gives $k_1 = m - (t - 1)/2$, so $j = (t - 1)/2$.

If $j \ge k_0(x)$, the argument in the proof of Lemma~\ref{lem:onset-three} gives $q < Q_*(x)$.
There $Q_*(x)$ decreases with $x$ and $\log_2 Q_*(1) < 1651$, so $q < 2^{1651}$.
If $j < k_0(x)$, then $j \le 518{,}368{,}290$.
If $\lvert y - x\, 2^{j}\, 2^{1/2} \rvert > 1$, the same argument gives
\[
q < 3\, x\, 2^{j} \le 3 \cdot 2^{17{,}631{,}663 + 518{,}368{,}290} = 3 \cdot 2^{535{,}999{,}953} < 2^{535{,}999{,}955}.
\]
If $\lvert y - x\, 2^{j}\, 2^{1/2} \rvert \le 1$, it gives
\[
\log_2 q < 2\log_2 x + 2j + 4 \le 2 \cdot 17{,}631{,}663 + 2 \cdot 518{,}368{,}290 + 4 = 1{,}071{,}999{,}910.
\]
In every case $q < 2^{1{,}071{,}999{,}971}$, which contradicts $q \ge 2^{1{,}071{,}999{,}971} + 1$.
Therefore $x > 2^{17{,}631{,}663}$.
\end{proof}

\begin{named}{Lemma 14.}
\marklemma{14}{lem:linear-two}
Let the cycle be other than $1 \to 3$, $35 \to 99$, and $27 \to 611$, and write $h = 2m + 1$.
Then
\[
D > 2^{0.535 m + 1},
\qquad
t > 0.535 m + \log_2 3,
\qquad
k_1 < \frac{1.465}{2}\, m - \frac{\log_2 3 - 1}{2},
\qquad
x < 2^{0.465 m + 1/2}.
\]
Also $t \ge 573{,}519{,}987$, $k_1 \le \lfloor \log_2 q \rfloor - 286{,}759{,}993$, and $k_2 \ge \lfloor \log_2 q \rfloor + 286{,}759{,}994$.
\end{named}

\begin{proof}
Proposition~\ref{prop:other-two} gives $D > 2^{0.535 m + 1}$ and $t > \log_2(3q/x)$.
Lemma~\ref{lem:length-2} gives $x = (q + 2^{k_1})/D$ and $2^{k_1} \le q - 1$.
Hence $q + 2^{k_1} \le 2q - 1 < 2q$, so $x < 2q/D$ and $3q/x > 3D/2$.
Therefore
\[
t > \log_2(3q/x) > \log_2(3D/2) = \log_2 D + \log_2 3 - 1.
\]
Since $\log_2 D > 0.535 m + 1$,
\[
t > 0.535 m + \log_2 3.
\]
Lemma~\ref{lem:length-2} gives $h = 2k_1 + t$, so $k_1 = (2m + 1 - t)/2 = m - (t - 1)/2$.
Thus
\begin{align*}
k_1
&< m - \frac{0.535 m + \log_2 3 - 1}{2} \\
&= \frac{1.465}{2}\, m - \frac{\log_2 3 - 1}{2}.
\end{align*}
Since $q = \lfloor 2^{h/2} \rfloor = \lfloor 2^m \sqrt{2} \rfloor$ and $2^m \sqrt{2}$ is irrational, $q < 2^m \sqrt{2}$ and $\log_2 q < m + 1/2$.
Combined with $x < 2q/D$ and $D > 2^{0.535 m + 1}$,
\[
\log_2 x < 1 + \log_2 q - \log_2 D < 1 + \bigl(m + \tfrac12\bigr) - \bigl(0.535 m + 1\bigr) = 0.465 m + \tfrac12,
\]
and $x < 2^{0.465 m + 1/2}$.

Proposition~\ref{prop:scan-x} gives $m \ge 1{,}071{,}999{,}971$.
At that value, $0.535 m + 1 = 573{,}519{,}985.485$, so $D > 2^{573{,}519{,}985}$.
Lemma~\ref{lem:length-2} gives $D \le 2^t - 1$, so $t > 573{,}519{,}985$.
Since $t$ is odd, $t \ge 573{,}519{,}987$.
For every larger $m$ the same lower bound on $D$ is larger.
Since $t$ is odd, $(t - 1)/2$ is an integer, and $k_1 = m - (t - 1)/2 \le m - 286{,}759{,}993$.
As in the proof of Proposition~\ref{prop:other-two}, $\lfloor \log_2 q \rfloor = m$, so $k_1 \le \lfloor \log_2 q \rfloor - 286{,}759{,}993$.
Also $k_2 = m + (t + 1)/2 \ge m + 286{,}759{,}994 = \lfloor \log_2 q \rfloor + 286{,}759{,}994$.
\end{proof}

\section{Longer cycles}
\label{sec:longer}

Length $3$ is the longest length for which every cycle through a constant odd $x$ falls under a root controlled by Lemma~\ref{lem:fixed-odd} or by Corollary~1.4 of~\cite{bauer-bennett-2002}.
Two thinner classes go further.
They depend on $\ell$ and on $h$ modulo $\ell$, so they are a corollary rather than part of Theorem~\ref{thm:bound}.

\begin{named}{Corollary 15.}
\marklemma{15}{cor:longer}
Let $x_0$ be an odd positive integer.
There is an integer $Q(x_0, \ell)$, read off from $x_0$ and $\ell$ by the comparison in the proof of Theorem~\ref{thm:bound}, with the following property.
If $x_0$ lies on a cycle of length $\ell$ with halving total $h$, and either
\begin{enumerate}
\item $\ell = 2^a$ with $a \ge 2$ and $h \equiv \ell/2 \pmod{\ell}$, or
\item $\ell = 3 \cdot 2^a$ with $a \ge 1$ and $h \equiv \ell/3$ or $2\ell/3 \pmod{\ell}$,
\end{enumerate}
then $q < Q(x_0, \ell)$.
\end{named}

\begin{proof}
In the first family every irrational root that occurs is $\sqrt{2}$.
An element aimed at an integer power of $2$ is large because it is odd.
Corollary~1.4 of~\cite{bauer-bennett-2002} supplies the exponent $\mu_2 = 1.465 < \varphi$ with odd part $x_0$.
The comparison in the proof of Theorem~\ref{thm:bound} then produces a finite check of the exponents below $j_0(x_0)$ and an explicit conflict for the larger exponents.

In the second family the only irrational roots that occur are $2^{1/3}$ and $2^{2/3}$.
Length $3$ is the case $a = 0$, in which those two residues are every nonzero class modulo $3$, and the restriction on $h$ disappears.
For $a \ge 1$, Lemma~\ref{lem:fixed-odd} supplies the same exponent $\mu_3 < \varphi$ with odd part $x_0$ or $3x_0$.
The comparison in the proof of Theorem~\ref{thm:bound} again produces the integer $Q(x_0, \ell)$.
\end{proof}

Length $4$ with $h$ odd is not included.
The first family uses Corollary~1.4 of~\cite{bauer-bennett-2002}.
The second uses Lemma~\ref{lem:fixed-odd}.

Corollaries~12, 13 and~14 of~\cite{hackman-i} are the case $x_0 = 1$.
Corollary~12 of that paper, and Theorem~1 of that paper for length $3$, conclude that no such $q$ exists.
Corollary~\ref{cor:longer} gives a bound in terms of $x_0$.
For $x_0 = 1$ the companion paper already carries out the search that rules out a length-$3$ cycle.
For a general $x_0$ the bound $Q(x_0, \ell)$ does not by itself rule the cycle out.
Corollaries~13 and~14 of~\cite{hackman-i} keep the constants of those papers ($c = 1$ and $j_0 = 9$ for $\sqrt{2}$; $j_0 = 11000$ for the cube roots).
The integer $Q(x_0, \ell)$ still depends on $x_0$, so it gives no uniform bound on every cycle of that length and residue.

\section{A heuristic}
\label{sec:all-cycles}

A cycle of length $\ell \ge 2$ with $h$ halvings has to satisfy several conditions at once.
As above, $q^\ell < 2^h < (q+1)^\ell$, so $\ell$ does not divide $h$ and $q = \lfloor 2^{h/\ell}\rfloor$ is a single odd integer at least $3$.
Write $D = 2^h - q^\ell$.
Starting from any element $x$ of the cycle, the halving counts are a composition of $h$ into $\ell$ positive parts, and the partial sums $S_j$ begin at $x$.
Unfolding the cycle relations gives
\[
x = \frac{N}{D}, \qquad N := q^{\ell-1} + \sum_{j=1}^{\ell-1} q^{\ell-1-j}\, 2^{S_j}.
\]
Thus $D$ divides $N$, and $x$ is a positive odd integer.
Tracing from $x$ recovers those counts exactly, and the $\ell$ elements are odd and distinct.

If, for a random composition, $N$ is uniformly distributed modulo $D$, then $D$ divides $N$ with probability $1/D$.
Counting cycles up to rotation, the expected number at this pair $(h,\ell)$ is then about
\[
\frac{1}{\ell D}\binom{h-1}{\ell-1}.
\]
This is only a heuristic: $N$ is a sparse combination of powers of $2$ with coefficients powers of $q$, so the residue $N \bmod D$ can be constrained, and the remaining conditions only make a cycle rarer.

The estimate is large only when $D$ is small, that is when $q^\ell$ is extraordinarily close to a power of $2$.
The known cycles are exactly those coincidences:
\[
5^2 = 2^5 - 7, \qquad 5^3 = 2^7 - 3, \qquad 181^2 = 2^{15} - 7.
\]
The corresponding expectations are $4/(2\cdot 7) \approx 0.29$, $15/(3\cdot 3) = 5/3$, and $14/(2\cdot 7) = 1$, against one, two, and two cycles found.
The only consecutive perfect powers are $8$ and $9$, of the other sign, so $D \ge 3$ for $\ell \ge 2$, and $D = 3$ and $D = 7$ are as small as occur here.

The next-smallest values of $D$ do not produce cycles.
For $q = 11$ and $h = 7$ one has $11^2 = 2^7 - 7$, the same $D$ as for $q = 5$ and $q = 181$, and the estimate is $6/(2\cdot 7) \approx 0.43$; but $11 + 2^{k_1} \not\equiv 0 \pmod{7}$ for every positive $k_1 < 7$, so $D$ never divides $N$.
For $q = 45$ and $h = 11$, likewise $45^2 = 2^{11} - 23$ and $45 + 2^{k_1} \equiv 0 \pmod{23}$ forces $2^{k_1} \equiv 1 \pmod{23}$, hence $11 \mid k_1$, which is impossible for $0 < k_1 < 11$.
For $q = 3$ and $h = 5$, $3^3 = 2^5 - 5$ and the estimate is $0.4$; the length-$3$ search in Appendix~\ref{app:computations} finds no cycle.
For $q = 3$ and $h = 8$, $3^5 = 2^8 - 13$ and the estimate is $0.54$; the length-$5$ search there finds no cycle.
For $q = 19$ and $h = 17$, $19^4 = 2^{17} - 751$, the fourth-power near miss recorded in the companion paper~\cite{hackman-i}, and the estimate is $0.19$; the length-$4$ search there finds no cycle.

Summing the same estimate over every admissible pair with $2 \le \ell \le 20$ and $h \le 200$ gives a total of about $7$ expected cycles, of which about $3$ sit in the three known slots and the remaining $4$ are spread over pairs that Appendix~\ref{app:computations} and the companion paper~\cite{hackman-i} have already emptied.
For large $h$ a typical $D$ is of size $\ell \delta\, 2^{h(1-1/\ell)}$, with $\delta$ the fractional part of $2^{h/\ell}$, and $\binom{h-1}{\ell-1}/D$ decays exponentially in $h$.
Length $1$ is different: $D = 2^h - q = 1$ holds for every Mersenne number, and $N = 1$, so the division succeeds and produces $1 \to 1$.

For a cycle at large $q$ there is a further obstruction.
A random odd input has $E[v_2(qn+1)] = 2$, while a cycle needs mean $h/\ell \approx \log_2 q$.
For $q = 5$ this is $2.32$, a mild excess, and three cycles exist.
For $q = 181$ it is $7.5$, which occurs only because $D = 7$.
For larger $q$ the same mean requires a still smaller relative gap, which is the Diophantine condition already used in Theorem~2 of the companion paper and Theorem~\ref{thm:finite}.

\section{Open problems}
\label{sec:open}

\begin{enumerate}
\item \textbf{All cycles.}
Conjecture~\ref{conj:all} is open.
Theorem~\ref{thm:finite} gives, for each $\ell \ge 2$, only finitely many $q$ with a cycle of length $\ell$, with no bound and no classification.
A bound in terms of a constant odd element, for lengths $2$ and $3$, is Theorem~\ref{thm:bound}, and for the classes in Corollary~\ref{cor:longer} it is that corollary.
A uniform bound on $q$ remains open.
Appendix~\ref{app:denominator} records the constraints on $D = 2^h - q^\ell$ for an arbitrary cycle.
Appendix~\ref{app:narrow} records three further restrictions.
Each one leaves that uniform bound open, and leaves open the list of lengths at a single $q$.
Proposition~B.2 gives, for each fixed integer $V$, only finitely many cycles with $\ell/\gcd(h,\ell) \le V$.
It gives no bound.
A cycle with $\gcd(h,\ell) = 1$ has $\ell/\gcd(h,\ell) = \ell$, and the proposition does not bound those lengths.
Proposition~\ref{prop:exponent} produces that bound from a restricted irrationality exponent at most $1.3247$ for every root $2^{r/\ell}$ and every odd positive integer $x$.
For length $2$, Section~\ref{sec:length-2} leaves the case $u \ge 3$, with $q \ge 2^{1{,}071{,}999{,}971}$ and every odd element greater than $2^{17{,}631{,}663}$.
The heuristic of Section~\ref{sec:all-cycles} is not a proof, and the searches leave infinitely many pairs $(\ell, h)$ unchecked.
\item \textbf{Novelty.}
Franco's 1990 thesis~\cite{franco-thesis} on Diophantine approximation and the $qx+1$ problem could not be consulted and should be checked for overlap, including with Theorem~\ref{thm:finite}.
\end{enumerate}

\appendix

\section{Restrictions on the denominator of any cycle}
\label{app:denominator}

Throughout this appendix, let the cycle have length $\ell \ge 2$, and let $D = 2^h - q^\ell$.
Lemma~\ref{lem:length-3-gap} gives a lower bound on $D$ when $\ell = 3$.
Lemma~\ref{lem:length-2} gives, when $\ell = 2$, that $D$ divides $2^t - 1$ and that $D \equiv 7 \pmod{8}$.
Proposition~A.3 is that divisibility for an arbitrary length.

\begin{named}{Proposition A.1.}
Then $h \ge 4$.
The integer $D$ is odd and $D \ge 3$.
Also $\gcd(D, q) = 1$.
$D \equiv 2^h \pmod{q}$ and $D \equiv -q^\ell \pmod{8}$.
In particular, if $\ell$ is even, then $D \equiv 7 \pmod{8}$.
If $\ell$ is even and $h$ is odd, then every prime divisor of $D$ is congruent to $1$ or $7$ modulo $8$.
\end{named}

\begin{proof}
Lemma~\ref{lem:any-window} gives $q^\ell < 2^h$.
Since $q \ge 3$ and $\ell \ge 2$, one has $2^h > 9$, so $h \ge 4$ and $D > 0$.
Thus $8$ divides $2^h$, and $D \equiv -q^\ell \pmod{8}$.
The integer $q$ is odd, so $q^2 \equiv 1 \pmod{8}$.
If $\ell$ is even, then $q^\ell = (q^2)^{\ell/2} \equiv 1 \pmod{8}$, and $D \equiv -1 \equiv 7 \pmod{8}$.
Here $2^h$ is even and $q^\ell$ is odd, so $D$ is odd.
Suppose $D = 1$.
Then $2^h - q^\ell = 1$, and both exponents are at least $2$.
Mihăilescu's theorem (2004)~\cite{mihailescu} says that the only solution of $a^u - b^v = 1$ in integers $a, b > 0$ and $u, v > 1$ is $3^2 - 2^3 = 1$.
The equation $2^h - q^\ell = 1$ would be a solution of the other sign.
Thus $D \ge 3$.
Also $\gcd(D, q) = \gcd(2^h, q)$.
Since $q$ is odd, this greatest common divisor equals $1$.
Reducing $D = 2^h - q^\ell$ gives $D \equiv 2^h \pmod{q}$.

Suppose $\ell$ is even and $h$ is odd.
Let $p$ be a prime divisor of $D$.
Then $q^\ell = 2^h - D \equiv 2^h \pmod{p}$.
Since $\gcd(D, q) = 1$, the prime $p$ does not divide $q$.
Since $D$ is odd, $p \ne 2$, and $p$ does not divide $2^h$.
The integer $\ell$ is even, so $q^\ell$ is a nonzero square modulo $p$, and therefore $2^h$ is a nonzero square modulo $p$.
Since $h$ is odd, $2$ is a quadratic residue modulo $p$, and therefore $p \equiv 1 \pmod{8}$ or $p \equiv 7 \pmod{8}$.
\end{proof}

Length $2$ has odd $h$, so every prime divisor of $D$ is congruent to $1$ or $7$ modulo $8$.
Lemma~\ref{lem:length-2} also gives that $D$ divides $2^t - 1$.

\begin{named}{Proposition A.2.}
Let $d = \gcd(h, \ell)$.
Let $X = 2^{h/d}$ and $Y = q^{\ell/d}$.
Then $D = X^d - Y^d$, so $X - Y$ divides $D$.
If $d = 1$, then $X - Y = D$.
If $d > 1$, then $X - Y > 1$ and $(X^d - Y^d)/(X - Y) > 1$.
\end{named}

\begin{proof}
The integers $h/d$ and $\ell/d$ are positive, and $X^d - Y^d = 2^h - q^\ell = D$.
The difference of powers factors as
\[
X^d - Y^d = (X - Y) \sum_{i=0}^{d-1} X^{d-1-i} Y^i.
\]
If $d = 1$, the sum equals $1$ and $X - Y = D$.

Now suppose $d > 1$.
Lemma~\ref{lem:any-window} gives that $\ell$ does not divide $h$, so $d < \ell$ and $\ell/d \ge 2$.
The same lemma gives $2^h > q^\ell$.
Since $q \ge 3$, one has $\log_2 q > 1$, so $h > \ell$ and $h/d > \ell/d \ge 2$.
The integer $h/d$ is therefore at least $3$, and both exponents $h/d$ and $\ell/d$ are at least $2$.
Also $X > Y$, so $X - Y$ is a positive integer.
Mihăilescu's theorem (2004)~\cite{mihailescu} says that the only solution of $a^u - b^v = 1$ in integers $a, b > 0$ and $u, v > 1$ is $3^2 - 2^3 = 1$.
The equation $2^{h/d} - q^{\ell/d} = 1$ would be a solution of the other sign.
Thus $X - Y > 1$.
The sum has $d \ge 2$ terms and its leading term is $X^{d-1}$.
Since $d - 1 \ge 1$ and $X = 2^{h/d} \ge 2^3$, that term is at least $8$, and the sum is greater than $1$.
\end{proof}

In particular, if $d > 1$, then $D$ is a product of two integers, each greater than $1$.
Each of the five cycles of length at least $2$ in the introduction has $\gcd(h, \ell) = 1$.

\begin{named}{Proposition A.3.}
Let $S_0 = 0$, and let
\[
f(z) = \sum_{j=0}^{\ell-1} 2^{S_j} z^{\ell-1-j}.
\]
Let $R$ be the resultant of $z^\ell - 2^h$ and $f(z)$.
Then $D$ divides $R$.
\end{named}

\begin{proof}
The relation $q x_{i-1} + 1 = 2^{k_i} x_i$ gives, by induction on $m$,
\[
2^{S_m} x_m = x_0 q^m + \sum_{j=0}^{m-1} 2^{S_j} q^{m-1-j}
\]
for $1 \le m \le \ell$.
For $m = \ell$ one has $x_\ell = x_0$ and $S_\ell = h$, so
\[
x_0(2^h - q^\ell) = \sum_{j=0}^{\ell-1} 2^{S_j} q^{\ell-1-j} = f(q).
\]
Thus $D$ divides $f(q)$.
The same identity $D = 2^h - q^\ell$ gives that $D$ divides $q^\ell - 2^h$.

The resultant $R$ is an integer.
The Sylvester matrix of $z^\ell - 2^h$ and $f(z)$ gives polynomials $u(z)$ and $v(z)$ with integer coefficients such that
\[
u(z)\, f(z) + v(z)\,(z^\ell - 2^h) = R.
\]
Substituting $z = q$ gives $R = u(q)\, f(q) + v(q)\,(q^\ell - 2^h)$.
Both terms on the right are divisible by $D$.

For $\ell = 2$ one has $f(z) = z + 2^{k_1}$, and $R = \pm(2^{2k_1} - 2^h)$.
Lemma~\ref{lem:length-2} gives $k_1 < k_2$.
Then $2^{2k_1} - 2^h = -2^{2k_1}(2^{k_2 - k_1} - 1)$.
Proposition~A.1 gives that $D$ is odd, so $D$ divides $2^{k_2 - k_1} - 1$.
\end{proof}

The last conclusion is the divisibility in Lemma~\ref{lem:length-2}.

\section{Finiteness past Theorem 2}
\label{app:narrow}

\begin{named}{Proposition B.1.}
Let $q$ be a Mersenne number, and write $q = 2^e - 1$ for an integer $e \ge 2$.
Let $\ell \ge 2$.
On a cycle of length $\ell$, $\ell > 1/\log_2(1+1/q)$ and $\ell/\gcd(h,\ell) > 1/\log_2(1+1/q)$.
\end{named}

\begin{proof}
Lemma~\ref{lem:any-window} gives $q^\ell < 2^h < (q+1)^\ell$.
Here $(q+1)^\ell = 2^{e\ell}$, so $h \le e\ell - 1$, and $h > \ell\log_2 q$.
Suppose $\ell\log_2(1+1/q) \le 1$.
Then
\[
\ell\log_2 q = e\ell - \ell\log_2\bigl(1 + \tfrac1q\bigr) \ge e\ell - 1,
\]
so $h > \ell\log_2 q \ge e\ell - 1 \ge h$, which is impossible.
Thus $\ell\log_2(1+1/q) > 1$.

Since $h \le e\ell - 1$ and $h > \ell\log_2 q$, the positive integer $s = e\ell - h$ satisfies
\[
1 \le s < \ell\log_2\bigl(1 + \tfrac1q\bigr).
\]
Let $v = \ell/\gcd(h,\ell)$.
Then $\gcd(h,\ell) = \gcd(s,\ell) \le s$, so $v$ satisfies
\[
v \ge \frac{\ell}{s} > \frac{1}{\log_2(1+1/q)}.
\]
\end{proof}

For $q = 3$ one has $2^{1/3} < 4/3 < \sqrt{2}$, because $2 < 64/27$ and $16/9 < 2$.
Thus $2 < 1/\log_2(4/3) < 3$, and the proposition excludes a cycle of length $2$.
It does not exclude length $3$.
The same lower bound is larger than $q\log 2$, because $\log(1+1/q) < 1/q$, so the excluded lengths grow with $q$.
Once $1/\log_2(1+1/q) \ge V$, a cycle for this $q$ has $\ell/\gcd(h,\ell) > V$, and Proposition~B.2 does not apply to it.

\begin{named}{Proposition B.2.}
Fix an integer $V \ge 2$.
Only finitely many cycles satisfy $\ell/\gcd(h,\ell) \le V$.
\end{named}

\begin{proof}
The odd elements of a cycle of length $\ell \ge 2$ are distinct odd positive integers, so
\[
\sigma \le \sum_{k=0}^{\ell-1} \frac{1}{2k+1}.
\]
The summand $f(t) = 1/(2t+1)$ is strictly decreasing for $t \ge 0$.
For each integer $k \ge 1$ and every $t \in [k-1,k)$, one has $f(t) > f(k)$, so
\[
\sum_{k=1}^{\ell-1} \frac{1}{2k+1} < \int_0^{\ell-1} \frac{dt}{2t+1} = \tfrac12\log(2\ell-1).
\]
The $k = 0$ term equals $1$, and therefore $\sigma < 1 + \tfrac12\log(2\ell-1)$.
Let $\Sigma$ be the positive real number $1 + \tfrac12\log(2\ell-1)$.
Lemma~\ref{lem:any-window} gives $(q+\delta)^\ell = \prod_i(q + 1/x_i)$, so
\[
\ell\log\Bigl(1 + \frac{\delta}{q}\Bigr) = \sum_i \log\Bigl(1 + \frac{1}{q x_i}\Bigr) < \frac{\sigma}{q}.
\]
For $u > 0$, $\log(1+u) > u/(1+u)$.
With $u = \delta/q$,
\[
\frac{\ell\delta}{q+\delta} < \frac{\sigma}{q},
\]
and therefore $\delta(\ell - \sigma/q) < \sigma$.
The series $e^u > u^2/2$ at $u = 2\sqrt{\ell}$ gives $e^{2\sqrt{\ell}} > 2\ell > 2\ell-1$, so $\log(2\ell-1) < 2\sqrt{\ell}$ and $\Sigma < 1 + \sqrt{\ell}$.
Also $\Sigma < \ell$.
Since $q \ge 3$, one has $\sigma < q\ell/2$, hence $\ell - \sigma/q > \ell/2$ and
\[
\delta < \frac{2\sigma}{\ell} < \frac{2\Sigma}{\ell}.
\]

Let $\varepsilon$ be the constant rational $1/10$.
Write $h = \ell m + r$ with integers $m \ge 0$ and $0 < r < \ell$, and set $\theta = 2^{r/\ell}$.
Let $v = \ell/\gcd(h,\ell)$.
The hypothesis is $v \le V$.
Let $g = \gcd(r,\ell)$.
Then $\gcd(h,\ell) = g$ and $\theta = 2^{\rho/v}$ with $\rho = r/g$, so $\gcd(\rho,v) = 1$ and $1 \le \rho \le v-1$.
In particular $1 < \theta < 2$ and $2^m\theta = q+\delta$, so $2^m = (q+\delta)/\theta < q+1$.
The integer $2^m$ is therefore at most $q$.
The hypothesis $v \le V$ leaves only finitely many such real numbers $\theta$.
Each is an algebraic irrational.
Let $\Phi$ be the set of real numbers $2^{\sigma_0/v_0}$ with integers $v_0,\sigma_0$ such that $2 \le v_0 \le V$ and $1 \le \sigma_0 \le v_0-1$.
That set is finite.
Each element lies strictly between $1$ and $2$.
A rational value would be an integer, so each element is an algebraic irrational.

The series for $e^{7/10}$ through the term of degree $4$ equals $482921/240000 > 2$, so $\log 2 < 7/10$.
The series for $\log(1+1/4)$ alternates, and the term after $(1/4)^3/3$ is negative, so $\log(5/4) < 1/4 - 1/32 + 1/192 = 43/192$.
Hence $\log 10 = 3\log 2 + \log(5/4) < 21/10 + 43/192 = 2231/960$.
The partial sum of $e$ through $1/8!$ equals $109601/40320$.
The cross product $109601 \cdot 500 - 40320 \cdot 1359 = 5620$ is positive, so $e > 2718/1000$ and $e^2 > 7387524/1000000$.
For $0 < z \le 1/25$, the remainder of $e^z$ after the quadratic term is at most $z^3/(6(1-z)) < z^2/2$, so $e^z < 1 + z + z^2$.
Thus $e^{1/25} < 651/625$ and $7e^{1/25} < 4557/625$.
The product $4557 \cdot 1600 = 7291200 < 7387524$, so $4557/625 < 7387524/1000000 < e^2$.
Hence $e^{49/25} > 7$ and $\log 7 < 49/25$.
Also $(6/5)^2 = 36/25$ and $(6/5)^4 = 1296/625 > 2$, so $(6/5)^8 > 4 > 3$ and $3^{1/10} < 6/5$.
The same fourth power gives $(6/5)^4 > 2$, so $2^{1/10} < 6/5$ and $2^{1/5} < 6/5$.

\topic{The length is short when $q < 4\ell$.}
Suppose $\delta \ge (q+1)^{-\varepsilon}$ and $q < 4\ell$.
Then $q+1 \le 4\ell$ and $\ell < 2\Sigma\,(4\ell)^{1/10}$.
Here $(4\ell)^{1/10} = 2^{1/5}\ell^{1/10} < (6/5)\ell^{1/10}$, and $\log(2\ell-1) < 7/10 + \log\ell$, so $\Sigma < 27/20 + \tfrac12\log\ell$.
Therefore
\[
2\Sigma\,(4\ell)^{1/10} < \frac{12}{5}\Bigl(\frac{27}{20} + \tfrac12\log\ell\Bigr)\ell^{1/10} = \Bigl(\frac{81}{25} + \frac65\log\ell\Bigr)\ell^{1/10}.
\]
The integer comparison $5^{10} = 9765625 > 7340032 = 7 \cdot 2^{20}$ is $(28/5)^{10} < 7^9$, so $7^{9/10} > 28/5$.
Combined with $\log 7 < 49/25$,
\[
7^{9/10} - \tfrac65\log 7 > \frac{28}{5} - \frac{6}{5}\cdot\frac{49}{25} = \frac{406}{125} > \frac{81}{25}.
\]
The derivative of $\ell^{9/10} - (6/5)\log\ell$ equals $(9\ell^{9/10} - 12)/(10\ell)$, which is positive for $\ell \ge 7$ because $7^{9/10} > 4/3$.
Thus $\ell^{9/10} - (6/5)\log\ell > 81/25$ for every real $\ell \ge 7$, and $2\Sigma\,(4\ell)^{1/10} < \ell$.
This contradicts $\ell < 2\Sigma\,(4\ell)^{1/10}$.
Hence $\ell \le 6$.
Then $q < 24$.
Lemma~\ref{lem:any-window} leaves finitely many integers $h$, and each such $h$ has finitely many compositions into $\ell$ positive parts.

\topic{A small fractional part.}
Suppose $\delta < (q+1)^{-\varepsilon}$.
Then $\delta < (2^m)^{-\varepsilon}$, because $2^m \le q < q+1$.
The fraction $q/2^m$ is in lowest terms, and
\[
\Bigl|\theta - \frac{q}{2^m}\Bigr| = \frac{\delta}{2^m} < (2^m)^{-1-\varepsilon}.
\]
Apply the $S$-unit form of Ridout's theorem (1958)~\cite{ridout}, with $S = \{2\}$ and this $\varepsilon$, to each element of the finite set of roots $2^{\rho/v}$.
Only finitely many reduced fractions meet the inequality.
For each such fraction the positive real number $\delta_* = 2^m\theta - q$ is positive, since $\theta$ is irrational.
A cycle producing that fraction has $\delta = \delta_*$, so $\ell < 2\Sigma/\delta_*$.
With $c = 2/\delta_*$ and $t = \sqrt{\ell}$, the bound $\Sigma < 1 + \sqrt{\ell}$ gives $t^2 < c(1+t)$, hence $t < c+1$ and $\ell < (c+1)^2$.
Each such fraction therefore allows only finitely many lengths.
Each pair $(q,\ell)$ again leaves finitely many compositions.

\topic{A large fractional part.}
Suppose now that $\delta \ge (q+1)^{-\varepsilon}$.
Then $\ell < 2\Sigma\,(q+1)^{\varepsilon}$.
The case $q < 4\ell$ has already been bounded, so assume $q \ge 4\ell$.
Lemma~\ref{lem:three-approx} applies.
The proof of that lemma records $e^{1/4} < 4/3$.
Since $\ell/q \le 1/4$ and $\delta \le e^{\ell/q}/a$ by Lemma~\ref{lem:any-window},
\[
a \le \frac{e^{\ell/q}}{\delta} \le e^{1/4}(q+1)^{\varepsilon} < \tfrac43(q+1)^{\varepsilon}.
\]
Let $y$ be the positive integer $\min(x_1,\ldots,x_{\ell-1})$.
Then $y > a$.

First let $q < 10^{300}$.
Then $q+1 \le 10^{300}$ and $(q+1)^{1/10} \le 10^{30}$, so $\ell < 2\Sigma \cdot 10^{30}$.
Let $\Sigma_*$ be the positive real number $1 + \tfrac12\log(2\ell)$.
Then $\Sigma < \Sigma_*$, and the derivative of $\ell/\Sigma_*$ equals $(\Sigma_* - 1/2)/\Sigma_*^2 > 0$, so $\ell/\Sigma_*$ is strictly increasing for $\ell > 1$.
At $\ell = 10^{32}$,
\[
\log(2\ell) = \log 2 + 32\log 10 < \tfrac{7}{10} + 32\cdot\frac{2231}{960} = \tfrac{7}{10} + \frac{2231}{30} < 76,
\]
so $\Sigma_* < 39$ and $10^{32}/\Sigma_* > 10^{32}/39 > 2 \cdot 10^{30}$.
For every $\ell \ge 10^{32}$ one therefore has $\ell > 2 \cdot 10^{30}\, \Sigma_* > 2 \cdot 10^{30}\, \Sigma$, a contradiction.
Thus $\ell < 10^{32}$.
Both $q$ and $\ell$ are bounded, and each pair leaves finitely many compositions.

Now let $q \ge 10^{300}$.
Then $q+1 < 2q$, so $(q+1)^{1/10} < (2q)^{1/10} < (6/5)q^{1/10}$ and $\ell < (12/5)\Sigma\, q^{1/10}$.
The series for $e$ is strictly less than $1 + 1 + \sum_{k \ge 1} 2^{-k} = 3$, so $e < 3$, $e^2 < 9 < 10$ and $\log 10 > 2$.
Suppose $\ell \ge q$.
Then $\Sigma < 1 + \sqrt{\ell}$ and $q \le \ell$ give $\ell < (12/5)(1+\sqrt{\ell})\,\ell^{1/10}$, so $\ell^{9/10} < (12/5)(1+\sqrt{\ell})$.
For $\ell \ge 1$ one has $1+\sqrt{\ell} \le 2\sqrt{\ell}$, hence $\ell^{2/5} < 24/5$ and $\ell < (24/5)^{5/2}$.
Here $(24/5)^2 = 576/25 < 24$ and $24^{1/2} < 5$, so $(24/5)^{5/2} < 120$.
This contradicts $\ell \ge 10^{300}$.
Thus $\ell < q$.
Then $\log(2\ell-1) < 7/10 + \log q$ and $\Sigma < 27/20 + \tfrac12\log q < \log q$, because $\log q > 600 > 27/10$.
Hence $\ell < (12/5)(\log q)\, q^{1/10}$.

Let $i$ be an index with $x_i = y$, and write $ir = \ell t + s$ with integers $t \ge 0$ and $0 \le s < \ell$.
Let $\varphi$ be the positive real number $2^{s/\ell}$, and let $F_i$ be the rational of Lemma~\ref{lem:three-approx}.
Set $G = F_i/2^t$.
Then $\theta_i = 2^t\varphi$ and $|\varphi - G| \le 2i\varphi/(aq)$.

Suppose $y \le q^{4/5}$ and $s = 0$.
Then $\varphi = 1$ and $\theta_i = 2^t$.
Let $b = S_i - im$ and $e = \max(im - S_i,\, 0)$.
Lemma~\ref{lem:three-approx} and $q \ge 4\ell > 4i$ give $|\theta_i - F_i| < \theta_i/2$, so $2^{t-1} < F_i < (3/2)2^t$.
The rational $F_i$ is not equal to $2^t$: otherwise $y = a$, contradicting distinctness.
If $b > t$, then $F_i = y\, 2^b/a$ and $|a - y\, 2^{b-t}| \le 2i/q < 1/2$.
The integer $y\, 2^{b-t}$ is even and $a$ is odd, so those two integers differ by at least $1$, which is impossible.
If $0 \le b \le t$, the same multiplication gives $|a\, 2^{t-b} - y| \le 2i\, 2^{t-b}/q$.
The left side is a positive integer, so $q \le 2i\, 2^{t-b}$.
From $F_i > 2^{t-1}$ one has $2^{t-b} < 2y/a$, hence $q < 4iy \le 4\ell\, q^{4/5}$ and $q^{1/5} < 4\ell$.
Then $q^{1/10} < (48/5)\log q$.
For $q \ge 10^{300}$ the left side is $10^{30}$ and the right side is less than $10 \cdot 300 \cdot 2231/960 < 10^4$, a contradiction.
The case $e \ge 1$ is the same comparison: Lemma~\ref{lem:three-approx} supplies $2^e \ge q/(2i\theta_i)$ and $F_i > \theta_i/2$, so $y > aq/(4i)$ and again $q < 4\ell\, q^{4/5}$.

Suppose $y \le q^{4/5}$ and $s \ge 1$.
Then $\varphi > 1$.
Also $\varphi = 2^{\sigma_0/v}$ for the integer $\sigma_0 = (i\rho \bmod v)$, so $\varphi \in \Phi$.
Since $q \ge 4\ell > 4i$ and $a \ge 1$, one has $|\varphi - G| < \varphi/2$, hence $G > \varphi/2 > 1/2$ and $G < 3\varphi/2 < 3$.
Let $E = im - S_i + t$, so $G = y/(a\, 2^E)$.
Write $G = P/d$ in lowest terms.
The odd part of $d$ divides $a$, so $d \cdot |d|_2 \le a$.
With $S = \{2\}$, the left side of the general form of Ridout's theorem is
\[
d^2\, |\varphi - G| \cdot |d|_2 \le d \cdot a \cdot \frac{4\ell}{aq} = \frac{4\ell\, d}{q},
\]
where $2i\varphi < 4\ell$.
If $E \ge 0$, then $a\, 2^E = y/G < 2y$, so $d < 2y$ and the left side is less than $8\ell y/q$.
If $E < 0$, then $d \le a < y$ and the left side is again less than $8\ell y/q$.
In both cases $\max(|P|,\, d) < 3y$: the numerator is $y$ when $E \ge 0$, and when $E < 0$ the identity $y\, 2^{-E} = Ga < 3a$ bounds the numerator by $3a < 3y$.
The inequality $8\ell y/q < (3y)^{-\varepsilon}$ therefore puts $G$ under the general form.
It rearranges as $8\ell\, 3^{1/10}\, y^{11/10} < q$.
Since $y \le q^{4/5}$ and $(4/5)(11/10) = 22/25$, the most demanding value is the right endpoint, and it is enough that $8\ell\, 3^{1/10} < q^{3/25}$.
Here $3^{1/10} < 6/5$ and $\ell < (12/5)(\log q)\, q^{1/10}$, so the left side is less than $(576/25)(\log q)\, q^{1/10}$.
The bound $\log q < 300 \cdot 2231/960 < 700$ and $q^{1/50} = 10^6$ give
\[
\frac{576}{25}\log q < 24 \cdot 700 < 10^6 = q^{1/50}.
\]
Thus $8\ell y/q < (3y)^{-\varepsilon} < \max(|P|,\, d)^{-\varepsilon}$.
Apply the general form of Ridout's theorem (1958)~\cite{ridout}, with $S = \{2\}$ and $\varepsilon = 1/10$, to each element of $\Phi$.
Only finitely many reduced fractions meet it.
For each such fraction $|\varphi - P/d|$ is a positive real number, because $\varphi$ is irrational.
The bound $|\varphi - G| \le 4\ell/(aq)$ then gives $q$ at most a constant multiple of $\ell$.
Combined with $\ell < (12/5)(\log q)\, q^{1/10}$, the real number $q^{9/10}/\log q$ is bounded, and so is $q$.
Each of those finitely many $q$ leaves finitely many lengths and finitely many compositions.

Suppose finally that $y > q^{4/5}$.
Then $1/y < q^{-4/5}$.
Also $2\ell < q^{1/5}$, because $(24/5)\log q < q^{1/10}$ for $q \ge 10^{300}$, so $2\ell/q < q^{-4/5}$ and
\[
\bigl|\theta - W\bigr| \le \frac{2}{q}\Bigl(\frac{1}{y} + \frac{2\ell}{q}\Bigr) < 4q^{-9/5},
\]
where $W = (\ell aq + 1)/(\ell a \cdot 2^m)$ is the rational of Lemma~\ref{lem:three-approx}.
The numerator $\ell aq + 1$ is congruent to $1$ modulo $\ell a$, so in lowest terms $W = P/d$ with $d = \ell a \cdot 2^e$ and $e = \max(m - v_2(\ell aq + 1),\, 0)$.
With $S = \{2\}$, the left side of the general form equals $(\ell a)^2 2^e\, |\theta - W| \le (\ell a)^2 q\, |\theta - W|$.
From $a < (4/3)(q+1)^{1/10} < (8/5)q^{1/10}$ and $\ell < (12/5)(\log q)\, q^{1/10}$,
\[
(\ell a)^2 < \frac{9216}{625}(\log q)^2\, q^{2/5}.
\]
Multiplying by $4q^{-4/5}$ gives a factor $36864/625 < 59$, so
\[
(\ell a)^2 q \cdot 4q^{-9/5} < 59\,(\log q)^2\, q^{-2/5}.
\]
Since $q+1 < 2q$,
\[
\max(|P|,\, d) < \ell a \cdot 2q + 1 < \frac{192}{25}(\log q)\, q^{6/5} + 1 < 16(\log q)\, q^{6/5}.
\]
The right side of the general form is therefore greater than $(16\log q)^{-1/10} q^{-3/25}$.
Comparing the two sides, it is enough that $59 \cdot 2 \cdot (\log q)^{21/10} < q^{7/25}$, because $16^{1/10} < 2$ and $2/5 - 3/25 = 7/25$.
For $q \ge 10^{300}$ one has $\log q < 700$, so $(\log q)^{21/10} < 700^3 = 343000000 < 10^9$, and $q^{7/25} = 10^{84}$.
Thus $59 \cdot 2 \cdot (\log q)^{21/10} < 10^{84}$.
Apply the general form of Ridout's theorem (1958)~\cite{ridout}, with $S = \{2\}$ and $\varepsilon = 1/10$, to each root $\theta = 2^{\rho/v}$.
Only finitely many reduced fractions meet it.
For each such fraction $|\theta - W|$ is a positive real number.
Lemma~\ref{lem:three-approx} and $y \ge 1$ give $|\theta - W| \le 2(1+2\ell)/q$, so $q$ is at most a constant multiple of $\ell$.
Combined with $\ell < (12/5)(\log q)\, q^{1/10}$, the real number $q^{9/10}/\log q$ is bounded, and so is $q$.
Each of those finitely many $q$ leaves finitely many lengths and finitely many compositions.
\end{proof}

The appeal to Ridout's theorem (1958)~\cite{ridout} gives no bound on $q$.
Since $\gcd(h,\ell) \ge 1$, every cycle of length at most $V$ satisfies $\ell/\gcd(h,\ell) \le V$.
Proposition~B.2 includes those cycles, and it includes a cycle of length greater than $V$ whenever $\gcd(h,\ell) \ge \ell/V$.
The hypothesis still requires $\ell/\gcd(h,\ell)$ to be at most a fixed integer.
Proposition~B.1 gives $\ell/\gcd(h,\ell) > 1/\log_2(1+1/q)$ whenever $q$ is a Mersenne number, so no single integer bounds $\ell/\gcd(h,\ell)$ for every such $q$.
If $\gcd(h,\ell) = 1$, then $\ell/\gcd(h,\ell) = \ell$.
The proposition does not bound those lengths.
In particular it does not bound the length of a cycle at $q = 3$.
At each fixed $q$, Theorem~\ref{thm:finite} still supplies finitely many cycles of each fixed length, and the argument above supplies finitely many cycles with $\ell/\gcd(h,\ell) \le V$.

\begin{named}{Proposition B.3.}
Let $S$ be a finite set of primes containing $2$.
Only finitely many cycles of length $\ell \ge 2$ have every prime factor of $q(2^h - q^\ell)$ in $S$.
\end{named}

\begin{proof}
Let $D = 2^h - q^\ell$.
Lemma~\ref{lem:any-window} gives $D > 0$.
Every prime factor of $qD$ lies in $S$, and $2 \in S$, so the positive rational numbers $q^\ell/2^h$ and $D/2^h$ are $S$-units.
Their sum is $1$.
Evertse (1984)~\cite{evertse} proves that $x + y = 1$ has at most $3 \times 7^{d+2s}$ solutions in $S$-units of a number field of degree $d$.
In that bound, $s$ is the number of places, including the infinite places.
Let $r$ be the number of primes in our set $S$.
Here the field is the field of rational numbers, so $d = 1$, and the places are the infinite place together with those $r$ primes, so $s = r+1$.
The number of solutions is at most $3 \times 7^{2r+3}$.
Each cycle in the proposition produces one of those solutions.
For each solution the integers $2^h$ and $q^\ell$ are determined.
A constant positive integer has finitely many representations as $q^\ell$ with $\ell \ge 2$ and $q$ an odd positive integer, since $\ell$ must divide every exponent in the prime factorization of that integer.
Each pair $(q,\ell)$ leaves finitely many compositions of $h$ into $\ell$ positive parts.
\end{proof}

The five cycles of length at least $2$ in the introduction have $q(2^h - q^\ell)$ divisible only by primes in $\{2,3,5,7,181\}$.
For that set one has $r = 5$, and the bound of Evertse (1984)~\cite{evertse} is $3 \times 7^{13} = 290667031221$.
The bound counts solutions of $x + y = 1$.
Several cycles may produce the same solution.
A cycle for which $q$ or $2^h - q^\ell$ has any other prime factor is not included.
For length $2$, Lemma~\ref{lem:length-2} gives that $D$ divides $2^t - 1$, so the prime factors of $D$ need not lie in a constant set.
A Mersenne number $q = 2^e - 1$ has $q + 1 = 2^e$ at length $1$, and those values of $q$ do not all have their prime factors in one constant finite set.

\section{Computations}
\label{app:computations}

Lemma~\ref{lem:any-window} gives $q^\ell < 2^h < (q+1)^\ell$ for a cycle of length $\ell \ge 2$, so $\ell \nmid h$ and $q = \lfloor 2^{h/\ell}\rfloor$.
A search of every $h \le H$ with $\ell \nmid h$ therefore excludes every such cycle with $q < 2^{\lfloor H/\ell \rfloor}$: it would have $2^h < (q+1)^\ell \le 2^{\ell \lfloor H/\ell \rfloor} \le 2^H$, so $h \le H$.
The time in each table is the search itself.

\topic{Any cycle.}
For each $\ell$ from $2$ to $9$, every $h$ up to the value in the table was tested, for odd elements that need not include $1$.
The only cycles found are those of Conjecture~\ref{conj:all}.
\begin{center}
\begin{tabular}{cccc}
\multicolumn{4}{c}{\textbf{Any cycle: only those of Conjecture~\ref{conj:all}}} \\[0.4em]
length $\ell$ & $h$ at most & for $q$ below & time \\
$2$ & $320{,}000$ & $2^{160{,}000}$ & $1{,}360$\,s \\
$3$ & $7{,}000$ & $2^{2{,}333}$ & $158$\,s \\
$4$ & $5{,}500$ & $2^{1{,}375}$ & $140$\,s \\
$5$ & $4{,}500$ & $2^{900}$ & $136$\,s \\
$6$ & $1{,}400$ & $2^{233}$ & $147$\,s \\
$7$ & $950$ & $2^{135}$ & $167$\,s \\
$8$ & $470$ & $2^{58}$ & $179$\,s \\
$9$ & $320$ & $2^{35}$ & $138$\,s \\
\end{tabular}
\end{center}
Length $2$ finds $1 \to 3 \to 1$ at $q = 5$, and $27 \to 611 \to 27$ and $35 \to 99 \to 35$ at $q = 181$.
Length $3$ finds $13 \to 33 \to 83 \to 13$ and $17 \to 43 \to 27 \to 17$ at $q = 5$.
Lengths $4$ through $9$ find no cycle.

\topic{A constant $q$.}
Here $q$ is constant and the length varies.
The only possible $h$ are the integers satisfying $q^\ell < 2^h < (q+1)^\ell$; if there is none, that length has no cycle.
For $q = 5$, lengths $2$ through $20$ produce only $1 \to 3 \to 1$ and the two cycles of length $3$ above.
For each odd $q$ from $7$ through $25$, lengths $2$ through $16$ produce no cycle.
The cycles at $q = 181$ were found in the length-$2$ row above, at $h = 15$, not in this search.
\begin{center}
\begin{tabular}{cccc}
\multicolumn{4}{c}{\textbf{Cycles at a constant $q$}} \\[0.4em]
$q$ & lengths & cycles found & time \\
$5$ & $2$ through $20$ & $1 \to 3 \to 1$;\ $13 \to 33 \to 83$;\ $17 \to 43 \to 27$ & under $1$\,s \\
$7$ & $2$ through $16$ & none & under $1$\,s \\
$9$ & $2$ through $16$ & none & under $1$\,s \\
$11$ & $2$ through $16$ & none & under $1$\,s \\
$13$ & $2$ through $16$ & none & under $1$\,s \\
$15$ & $2$ through $16$ & none & under $1$\,s \\
$17$ & $2$ through $16$ & none & under $1$\,s \\
$19$ & $2$ through $16$ & none & under $1$\,s \\
$21$ & $2$ through $16$ & none & under $1$\,s \\
$23$ & $2$ through $16$ & none & under $1$\,s \\
$25$ & $2$ through $16$ & none & under $1$\,s \\
\end{tabular}
\end{center}

\topic{The bounds $Q_2(x)$ and $Q_3(x)$.}
For every odd $x \le 35$, the values of $Q_{3,0}(x)$ were computed with upper limit $11000$ on $j$.
The gate $2^{11001}$ does not appear in $Q_3(x)$.
The quantity $Q_{3,*}(x)$ is smallest at $x = 3$, where $\log_2 Q_{3,*}(3) = 17906.510814$, so $Q_{3,*}(x) > 2^{17906} > 2^{11001}$ for every odd positive integer $x$.

\topic{Powers of two above $Q_{3,*}(x)$.}
The integer $\lceil Q_{3,*}(x) \rceil$ is the sharper bound, and it has thousands of digits.
The table uses the next power of $2$ above $Q_{3,*}(x)$, read off from the logarithm.
For $x = 5$ one has $\log_2 Q_{3,*}(5) = 18213.724287504393$, so $\lceil Q_{3,*}(5) \rceil$ has $5483$ digits and begins $753951296\ldots$, while the table entry is $2^{18214}$.
The same rounding, $2^{\lceil \log_2 \,\cdot\, \rceil}$, is used for $Q_{3,0}(x)$, $Q_{2,*}(x)$ and $Q_2(x)$.

For every odd $x$ from $1$ to $35$, the entry $2^{\lceil \log_2 Q_{3,*}(x) \rceil}$ is:

\begin{center}
\begin{tabular}{rrrrrr}
$x$ & & $x$ & & $x$ & \\
$1$ & $2^{17990}$ & $13$ & $2^{18387}$ & $25$ & $2^{18514}$ \\
$3$ & $2^{17907}$ & $15$ & $2^{18190}$ & $27$ & $2^{18304}$ \\
$5$ & $2^{18214}$ & $17$ & $2^{18439}$ & $29$ & $2^{18543}$ \\
$7$ & $2^{18273}$ & $19$ & $2^{18460}$ & $31$ & $2^{18556}$ \\
$9$ & $2^{18094}$ & $21$ & $2^{18255}$ & $33$ & $2^{18344}$ \\
$11$ & $2^{18356}$ & $23$ & $2^{18497}$ & $35$ & $2^{18580}$ \\
\end{tabular}
\end{center}

Rows in which $3$ divides $x$ sit below their neighbors because $s' = x/\gcd(x, 3)$.
The row $x = 3$ has $s' = 1$ and $c_3(x) = c_3(3)$.
The value $c_3(3)$ is the place this table uses Lemma~\ref{lem:fixed-odd} past odd part $1$ or $3$.

\topic{Powers of two above $Q_{3,0}(x)$.}
For the same odd $x$, with the maximum taken over $0 \le j < 11000$, the entry $2^{\lceil \log_2 Q_{3,0}(x) \rceil}$ is:

\begin{center}
\begin{tabular}{rrrrrr}
$x$ & & $x$ & & $x$ & \\
$1$ & $2^{11003}$ & $13$ & $2^{11007}$ & $25$ & $2^{11009}$ \\
$3$ & $2^{11002}$ & $15$ & $2^{11006}$ & $27$ & $2^{11006}$ \\
$5$ & $2^{11006}$ & $17$ & $2^{11009}$ & $29$ & $2^{11012}$ \\
$7$ & $2^{11006}$ & $19$ & $2^{11009}$ & $31$ & $2^{11008}$ \\
$9$ & $2^{11006}$ & $21$ & $2^{11006}$ & $33$ & $2^{11005}$ \\
$11$ & $2^{11005}$ & $23$ & $2^{11006}$ & $35$ & $2^{11008}$ \\
\end{tabular}
\end{center}

In every row the exponent $j$ that attains the maximum lies in $\{10997, 10998, 10999\}$.
For these eighteen values, $Q_{3,0}(x) < Q_{3,*}(x)$.
That comparison is checked for odd $x \le 35$.
It is not proved for every odd $x$.
From $t^3 - 2^r$, if the distance from $s'\, 2^{j+r/3}$ to an integer is positive, that distance is greater than $1/(Q^2 \cdot 3(2^{2/3}+1)^2)$ once the integer is within $1$ of $s'\, 2^{j+r/3}$, with $Q = s' 2^j$.
With $j \le 10999$ and $s' \le x$ this only yields $Q_{3,0}(x) < x^3\, 2^{33003}$.
At $x = 1$ the two logarithms, of this upper estimate and of $Q_{3,*}(1)$, are $33003$ and $17989$.
The estimate $x^3\, 2^{33003}$ stays above $Q_{3,*}(x)$ until $x$ is about $2^{67}$.
Past that size, $Q_{3,*}(x)$ grows like $x^{1/\gamma_3}$ with $1/\gamma_3 \approx 229$, while $x^3\, 2^{33003}$ grows like $x^3$.

\topic{The value $Q_3(1)$.}
Here $s' = 1$, so $J_3(1) = 11000$.
The maximum $Q_{3,0}(1)$ is less than $2^{11003}$, and $Q_{3,*}(1)$ is less than $2^{17990}$, so
\[
Q_3(1) = Q_{3,*}(1) < 2^{17990}.
\]
Each $h$ gives one candidate $q = \lfloor 2^{h/3} \rfloor$.
The bound $q < 2^{17990}$ means it is enough to check $h < 53970$.
The companion paper checks every $h \le 200{,}000$ and finds no cycle, so that search includes the range.
Theorem~1 of~\cite{hackman-i} is this case.

The proof of that theorem in the companion paper reaches the larger cutoff $q < 2^{22003}$.
There the symbol $x$ is the element after $1$.
The argument asks for that element to clear $2^{11001}$, using a lower bound of the shape $\sqrt{q/2}$.
The integer $Q_3(1)$ does not use that lower bound.
Once $q > Q_{3,0}(1)$, an exponent below $11000$ is already impossible, and the conflict is $Q_{3,*}(1)$.

\topic{Powers of two above $Q_{2,*}(x)$.}
This integer is used in the proof and is at most $Q_2(x)$.
As $x \to \infty$, $\log_2 Q_{2,*}(x)$ decreases to $\mu_2/\gamma_2 \approx 6.733$.
The row $x = 11$ has logarithm $7.013$, so that one ceiling moves if the fourth digit of the exponent $1.465$ is lower.
For odd $x$ from $1$ to $35$, the entry $2^{\lceil \log_2 Q_{2,*}(x) \rceil}$ is $2^{11}$ at $x = 1$, then $2^{8}$ for $x = 3, 5, 7, 9, 11$, and $2^{7}$ for every odd $x$ from $13$ to $35$.

\topic{Values of $Q_2(x)$.}
For these eighteen values the last failure of $|\sqrt{2} - p/(x\, 2^j)| > 2^{-1.465 j}$ lies at $j \le 21$, and a check through $j < 53620$ finds no later failure.
The entry is $2^{\lceil \log_2 Q_2(x) \rceil}$.

\begin{center}
\begin{tabular}{rrrr}
$x$ & $j_0(x)$ & winning $j$ & $Q_2(x)$ below \\
$1$ & $9$ & $7$ & $2^{15}$ \\
$3$ & $12$ & $11$ & $2^{18}$ \\
$5$ & $22$ & $17$ & $2^{33}$ \\
$7$ & $12$ & $10$ & $2^{18}$ \\
$9$ & $12$ & $11$ & $2^{18}$ \\
$11$ & $13$ & $12$ & $2^{19}$ \\
$13$ & $14$ & $9$ & $2^{21}$ \\
$15$ & $22$ & $17$ & $2^{32}$ \\
$17$ & $15$ & $13$ & $2^{23}$ \\
$19$ & $17$ & $14$ & $2^{25}$ \\
$21$ & $17$ & $12$ & $2^{25}$ \\
$23$ & $16$ & $14$ & $2^{23}$ \\
$25$ & $22$ & $17$ & $2^{32}$ \\
$27$ & $15$ & $14$ & $2^{22}$ \\
$29$ & $17$ & $16$ & $2^{25}$ \\
$31$ & $17$ & $16$ & $2^{25}$ \\
$33$ & $15$ & $13$ & $2^{22}$ \\
$35$ & $22$ & $17$ & $2^{32}$ \\
\end{tabular}
\end{center}

The cycle $27 \to 611 \to 27$ has $q = 181 < 2^{22}$.
The cycle $35 \to 99 \to 35$ has $q = 181 < 2^{32}$.

The powers of two above $Q_{3,*}(x)$ took under $1$\,s, and those above $Q_{2,*}(x)$ took under $1$\,s.
The powers of two above $Q_{3,0}(x)$, the maximum over $0 \le j < 11000$, took $43$\,s.
The table of $Q_2(x)$ took $47$\,s.

\topic{Digits of $\sqrt{2}$.}
Let $N = 1{,}072{,}000{,}000$ and let $I = \lfloor 2^N \sqrt{2} \rfloor$.
The integer $I$ was computed and checked by $I^2 < 2^{2N+1} < (I+1)^2$.
Place $1$ is the bit of $\sqrt{2}$ just below the integer bit.
For a cycle of length $2$ one has $h = 2m + 1$ and $2^{h/2} = q + \delta$, so the binary digits of $\delta$ are places $m+1$, $m+2$, $\ldots$ of $\sqrt{2}$.
Proposition~\ref{prop:other-two} covers every $m < 160{,}000$.
From $m = 160{,}000$ through $m = 1{,}071{,}999{,}970$, the longest run of consecutive zeros in those places has length $29$, at $m = 1{,}005{,}210{,}007$.
The scan took $134$\,s.

\begin{thebibliography}{99}
\bibitem{bauer-bennett-2002}
M.~Bauer and M.~A.~Bennett, Applications of the hypergeometric method to the generalized Ramanujan--Nagell equation, Ramanujan J.\ 6 (2002), 209--270, \doilink{10.1023/A:1015779301077}.

\bibitem{bauer-bennett-2018}
M.~Bauer and M.~A.~Bennett, Ramanujan--Nagell cubics, Rocky Mountain J.\ Math.\ 48 (2018), 385--412, \doilink{10.1216/rmj-2018-48-2-385}.

\bibitem{beukers}
F.~Beukers, On the generalized Ramanujan--Nagell equation~I, Acta Arith.\ 38 (1980/81), 389--410, \doilink{10.4064/aa-38-4-389-410}.

\bibitem{crandall}
R.~E.~Crandall, On the ``$3x+1$'' problem, Math.\ Comp.\ 32 (1978), 1281--1292, \doilink{10.1090/S0025-5718-1978-0480321-3}.

\bibitem{evertse}
J.~H.~Evertse, On equations in $S$-units and the Thue--Mahler equation, Invent.\ Math.\ 75 (1984), 561--584, \doilink{10.1007/BF01388644}.

\bibitem{franco-thesis}
Z.~Franco, Diophantine Approximation and the $qx+1$ Problem, PhD thesis, UC Berkeley, 1990.

\bibitem{hackman-i}
T.~Hackman, I. Cycles through 1 for the $qn+1$ map: no cycles of length $3$, and constraints on longer cycles, companion paper.

\bibitem{mihailescu}
P.~Mihăilescu, Primary cyclotomic units and a proof of Catalan's conjecture, J.~Reine Angew.~Math.~572 (2004), 167--195, \doilink{10.1515/crll.2004.048}.

\bibitem{ridout}
D.~Ridout, The $p$-adic generalization of the Thue--Siegel--Roth theorem, Mathematika 5 (1958), 40--48, \doilink{10.1112/S0025579300001339}.

\bibitem{santos}
R.~Santos, On the Collatz general problem $qn+1$, preprint (2020), \href{https://arxiv.org/abs/2005.00346}{\nolinkurl{arXiv:2005.00346}}.

\end{thebibliography}

\end{document}
